  \chapter{Thermal Effects on Breakdown and Transition}\label{chapter:breakdown}

Turbulence development in hypersonic boundary layers proceeds through a sequence beginning with linear instability growth and culminating in nonlinear breakdown. Whilst \gls{LST} captures the second Mack mode during early amplification, the subsequent nonlinear dynamics remain difficult to predict with comparable accuracy. Breakdown of coherent structures into turbulence determines the heating, frictional drag, and aerothermal performance of hypersonic vehicles. A characteristic feature is the overshoot in surface quantities such as skin-friction coefficient and Stanton number. Rather than the gradual increase predicted by linear theory, both exceed their fully turbulent values during transition before relaxing downstream \citep{Holden1972}. This behaviour arises from abrupt, localised enhancement of mixing and momentum transport during vortex breakup , which injects turbulent-like fluctuations into the near-wall region faster than the mean flow equilibrates. This transient amplification of aerodynamic heating poses a significant challenge for \gls{TPS} design and risks localised material failure. Understanding this response remains critical for hypersonic vehicle design across a wide range of flight regimes.

Second-mode instabilities are the dominant mechanism in high-Mach-number boundary layers. Controlled-forcing experiments using harmonic point sources confirmed the second mode as the primary instability and identified subharmonic resonance as a breakdown route \citep{Shiplyuk2003,Bountin2008}. More recent work in the Purdue quiet tunnel identified fundamental resonance as the dominant transition pathway on a Mach~6 flared cone \citep{Chynoweth2014} and established amplitude thresholds for turbulent spot formation \citep{Casper2011,Casper2012}. Low-enthalpy studies, including Mach~21 nitrogen-tunnel experiments \citep{Mironov2000}, revealed four-wave parametric resonance between fundamental and harmonic modes. However, such experiments remain limited by quiet-tunnel availability, perfect-gas conditions, and incomplete access to nonlinear dynamics. High-fidelity \gls{DNS} is necessary to systematically assess thermochemical nonequilibrium effects and wall thermal variations.

Temporal and spatial \gls{DNS} studies have investigated fundamental transition mechanisms in compressible hypersonic boundary layers. For instance, \citet{Mayer2011} computed Mach~3 transition triggered by oblique breakdown, a pathway first identified by \citet{Fasel1993}. Such simulations typically employ either broadband forcing across multiple frequencies or selective excitation of dominant energy modes \citep{AlamSandham2000,Bhaskaran2011}. Whilst the latter strategy does not replicate realistic flight or noisy wind tunnels, it efficiently isolates underlying physics. Fundamental and subharmonic resonance originate from amplification of a two-dimensional instability wave. Once this wave reaches finite amplitude, secondary oblique instabilities emerge, growing either at the same frequency (fundamental resonance) or at half the frequency (subharmonic resonance). The current study examines the transition mechanism through the second excitation method, building on \Cref{chapter:disturbance_growth}. Fundamental \gls{DNS} studies of turbulent boundary-layer behaviour in hypersonic flows were initiated by \citet{Martin2007} and advanced by \citet{Duan2011}, who examined the influence of wall temperature and Mach number on mean flow statistics in cold hypersonic flows. \citet{Franko2013} performed a spatial \gls{DNS} of a Mach~6 \gls{ZPG} boundary layer, examining the skin-friction and heat-transfer overshoot observed in experiments and identifying first-mode oblique breakdown as a plausible transition path. An extensive database of cooled boundary layers up to Mach~14 was later generated by \citet{Zhangetal2018}.

Fewer studies have incorporated real gas effects. \citet{Stemmer2005} simulated hypersonic boundary layer breakdown for reactive and non-reactive flow in a Mach~20 boundary layer. \citet{Marxenetal2013,Marxenetal2014} incorporated fundamental resonance and nonlinear modes in a Mach~10 adiabatic flat-plate flow, concluding that chemical reactions had no significant effect on growth rates. \citet{DiRenzo2021} performed a complete transition study incorporating chemical nonequilibrium and observed good qualitative agreement with previous works. \citet{Passiatoreetal2021} provided a comprehensive analysis of breakdown to turbulence under subharmonic second-mode excitation, identifying that fundamental mode saturation and energy transfer to oblique subharmonic modes lead to $\Lambda$-vortex formation. Whilst finite-rate chemistry had little effect on initial transition stages, the endothermic nature of later stages drained modal energy and delayed transition by approximately 7\%.

Several \gls{DNS} studies have examined real-gas effects on hypersonic boundary layers, though most focused on chemical nonequilibrium. \citet{DiRenzo2021} investigated transition in a Mach~10 hypersonic boundary layer at suborbital enthalpy, capturing the sequence of second-mode amplification, resonance with oblique harmonics, and breakdown, yielding four- to five-fold increases in wall friction and heating. That work also considered finite-rate chemistry in cooled turbulent boundary layers, whilst \citet{Passiatoreetal2021} extended the framework to pseudo-adiabatic walls to isolate chemical nonequilibrium effects. A significant limitation was the neglect of thermal nonequilibrium. More recently, \citet{Williamsetal2025} examined chemistry--turbulence coupling and showed that turbulence drives variations in chemical source terms whilst having minimal impact on turbulence structure. Most studies have been limited to fully thermochemical equilibrium conditions at unrealistically high Mach numbers, leaving understanding of isolated thermal effects incomplete.

This Chapter extends the \gls{LST} analysis to the breakdown stage across \gls{CPG}, \gls{TEQ}, and \gls{TNE} formulations. All cases are treated as chemically frozen, as the selected conditions do not trigger significant chemical activity \citep{BitterShepherd2015}. At high-enthalpy conditions, deviations from equilibrium become important, and it is necessary to assess how different thermal models affect breakdown onset and progression. Attention is given to the interplay between linear amplification, nonlinear resonance, secondary instability growth, and turbulent spot formation. Emphasis is placed on the overshoot behaviour of skin friction and heat transfer, the influence of cold-wall conditions in accelerating breakdown, and the collapse of Mack second modes into fully developed turbulence across the thermal sates. The analysis further examines the wall-temperature response of \gls{TNE} flows, where vibrational nonequilibrium modifies breakdown dynamics compared to \gls{CPG} and \gls{TEQ} cases. These results provide detailed a characterisation of the nonlinear transition regime under different thermal models, Mach numbers and wall conditions, highlighting the consequences for aerothermodynamic design and the challenges in developing predictive tools.

\section{Numerical setup and flow definitions}\label{section:breakdown:numericalsetup}


This study examines spatially developing \gls{ZPG} flat-plate boundary layers formed behind a shock wave of a wedge placed in a Mach 10 freestream at 36 km altitude, conditions representative of the heat-rate limit along a hypersonic glide-vehicle trajectory~\citep{Rizvi2015} as detailed in \Cref{section:disturbances:computationdomain}. Post-shock conditions of Mach 6 and 8 are investigated. The edge conditions are $p_e = 4.91~\mathrm{kPa}$ and $T_e = 620~\mathrm{K}$ for Mach 6, and $p_e = 1.78~\mathrm{kPa}$ and $T_e = 365~\mathrm{K}$ for Mach 8 as mentioned in~\Cref{table:flight_conditions}. The governing equations are the compressible Navier--Stokes equations under thermal nonequilibrium and chemically frozen assumptions. The computational domain is a rectangular box with streamwise extents of $800\delta^*_0$ (Mach 6) and $1600\delta^*_0$ (Mach 8), wall-normal height of $100\delta^*_0$, and spanwise width of $20\delta^*_0$. The grid is uniform in the streamwise and spanwise directions, with wall-normal spacing generated via hyperbolic tangent stretching, $y = L_y \sinh(s_1 \xi)/\sinh(s_1)$ for uniformly distributed $\xi \in [0,1]$ and stretching factor $s_1 = 5.0$. Grid resolutions are presented in \Cref{tab:domain_grid_dns} and \Cref{tab:domain_grid_iles}.

\begin{table}[hbt!]
\centering
\caption{Grid resolution and thermal conditions for DNS simulations.}
\begin{tabular}{lccccc}
\toprule
Case & Gas Model & $M_e$ & $T_w/T_e$ &  ($N_x \times N_y \times N_z$) & Total points \\
\midrule
\multirow{3}{*}{M6Tw18} & CPG & 6 & 2.90 & $5200 \times 600 \times 275$ & $8.58\times 10^8$ \\
 & TEQ & 6 & 2.90 & $5200 \times 600 \times 275$ & $8.58\times 10^8$ \\
 & TNE & 6 & 2.90 & $5200 \times 600 \times 275$ & $8.58\times 10^8$ \\
\bottomrule
\end{tabular}
\label{tab:domain_grid_dns}
\end{table}
\begin{table}[hbt!]
\centering
\caption{Grid resolution and thermal conditions for ILES simulations.}
\begin{tabular}{lccccc}
\toprule
Case & Gas Model & $M_e$ & $T_w/T_e$ &  ($N_x \times N_y \times N_z$) & Total points \\
\midrule
M6Tw12 & TNE & 6 & 1.93 & $2600 \times 600 \times 205$ & $3.20\times 10^8$ \\
M6Tw18 & TNE & 6 & 2.90 & $2300 \times 500 \times 165$ & $1.84\times 10^8$ \\
M8Tw12 & TNE & 8 & 3.29 & $4600 \times 500 \times 165$ & $3.68\times 10^8$ \\
M8Tw18 & TNE & 8 & 4.93 & $4600 \times 500 \times 165$ & $3.68\times 10^8$ \\
\bottomrule
\end{tabular}
\label{tab:domain_grid_iles}
\end{table}

Inflow boundary-layer profiles are obtained from a locally self-similar framework solved using Chebyshev collocation, ensuring consistent specification across \gls{CPG}, \gls{TEQ}, and \gls{TNE} conditions. For \gls{TNE}, a thermally frozen \gls{TFG} similarity solution provides the inflow as detailed in \Cref{section:similaritysolution}. Conservative variables are initialised and imposed via inlet pressure extrapolation. Outflow and far-field boundaries employ simple extrapolation, whilst spanwise boundaries are periodic. The wall is non-catalytic and isothermal at $T_w = 1200~\text{K}$ (cold) or $T_w = 1800~\text{K}$ (hot). Cold walls enhance second-mode amplification by thinning the boundary layer and destabilising the acoustic spectrum, whilst hot walls stabilise acoustic-trapped disturbances, delaying transition.

The inflow Reynolds number is $Re_{\delta^*_0} = 10{,}000$. The governing equations are discretised using a 5th-order WENO-Z scheme for convective terms and 4th-order central differences for viscous and metric terms, with 3rd-order SSP Runge--Kutta time integration. Transport properties follow the Musawi--Sandham viscosity model~\citep{Musawi2025} and vibrational relaxation uses the Millikan--White formulation with Park constants~\citep{Park1993}. Non-dimensional timesteps are $\Delta t^{*} = 0.012$ (Mach 6) and $\Delta t^{*} = 0.016$ (Mach 8) with a \gls{CFL} of $0.05$ across all cases, ensuring stability whilst resolving high-frequency acoustic waves characteristic of second-mode instability.

\subsection{Disturbance seeding and excitation method}

Transition to turbulence is induced via a suction--blowing forcing strip following \citep{AlamSandham2000}, extended here to a two-dimensional domain with large-amplitude perturbations. A wall mass-flux perturbation is imposed at $x_f$:
\begin{equation}
\rho v_w(x,z,t) \;=\; \rho a_0 \, g(x) \, h(z) \, \sum_{i=1}^{n} \sin(2\pi f_i t),
\end{equation}
where the streamwise envelope is
\begin{equation}
g(x) = \exp\left(-\frac{(x-x_f)^2}{2(\delta^*_0)^2}\right),
\end{equation}
and the spanwise modulation is
\begin{equation}
h(z) = 1 + 0.06 \left[\sin\!\left(\frac{4\pi z}{L_z}\right) + \sin\!\left(\frac{2\pi z}{L_z}\right)\right].
\end{equation}
The Gaussian envelope has width $2\sigma^2 = 4(\delta^*_0)^2$ and is centred at $x_f = 25\delta^*_0$ from the domain origin. The forcing amplitude is $a_0 = 0.1 \, u_e$, sufficiently large to trigger transition within the domain whilst remaining representative of realistic disturbance levels. The frequency set is case-dependent: for Mach 6, $f^{*}_i = \{0.15,\,0.16\}$, and for Mach 8, $f^{*}_i = \{0.10,\,0.11\}$. These frequencies are selected using linear stability theory to ensure second-mode growth reaches its maximum between $x/\delta^*_0 \approx 220$ and $400$ across all cases, coinciding with the forcing location at $x_f \approx 25\delta^*_0$. This alignment excites the most naturally amplified disturbances and sustains their growth within the computational domain.


\begin{figure}[t]
    \centering
    \includegraphics[width=9cm]{resources/figures/breakdown/insert_shapefunction_natural.pdf}
    \caption{Spatial structure of the wall mass-flux forcing function.}
    \label{figure:forcing_shapefunction}
\end{figure}

The spatial structure of the forcing function is shown in \Cref{figure:forcing_shapefunction} with a comparison to~\citep{Franko2013}. Rather than imposing oblique modes, three-dimensionality is introduced through spanwise modulation, which imposes long- and short-wavelength subharmonics ($\lambda_z = L_z/2$ and $L_z$). This approach is justified by the stability analysis in \Cref{figure:stability_plot}, which presents a Mach 6 case under different thermal excitations with adiabatic wall results in the left column and "hot" wall results in the right column. The analysis demonstrates the absence of first-mode instability across all thermal models; consequently, oblique modes arising from first-mode interactions need not be artificially introduced. Instead, $\Lambda$-shaped vortices and streaks emerge naturally, isolating the evolution of two-dimensional second-mode disturbances and their secondary instabilities without prescribing an artificial oblique structure.

\begin{figure}[t]
    \centering
    \includegraphics[width=12cm]{resources/figures/breakdown/result_mach6_stability_diagram.pdf}
    \caption{Comparison of Mach 6 stability diagrams at $\delta_{*}^{0}=255$ for (a) CPG, (b) TEQ, and (c) TNE models. For each model, the left column shows the adiabatic-wall case, while the right column corresponds to the isothermal 1800\,K wall condition.}
    \label{figure:stability_plot}
\end{figure}

Each simulation spans four consecutive flow-through times ($t_f = L_x/u_e$). The first two flow-throughs remove initial transients; the solution is restarted for two further flow-throughs whilst accumulating Favre-averaged statistics. This captures the complete sequence from controlled second-mode excitation through nonlinear interactions to turbulence onset.

\section{Influence of thermal excitations on flow behaviour}

The influence of thermal modelling on hypersonic boundary layer behaviour is investigated by comparing \gls{CPG}, \gls{TEQ}, and \gls{TNE} formulations for the Mach 6 hot-wall configuration. Linear stability analyses reveal that thermal modelling substantially affects second-mode growth rates. \gls{TEQ} exhibits the highest growth rates, exceeding those from \gls{TNE} and \gls{CPG}, due to stronger coupling between vibrationally equilibrated energy modes and acoustic disturbances. The \gls{TNE} formulation introduces vibrational relaxation effects that reduce modal amplification, whilst \gls{CPG} yields the weakest instability growth due to the absence of additional energy channels. These differences arise primarily from variations in the base flow. The resulting \gls{LST} predictions provide context for interpreting the nonlinear transition dynamics in the simulations.

Transition is characterised through skin-friction coefficient, Stanton number, mean velocity profiles, and velocity fluctuation components, with onset identified from the sharp rise in skin friction and Stanton number. Reynolds stress components and wall-normal heat flux are examined across the transitional and turbulent regions to assess how the three thermal models differ in turbulence production and identify the mechanisms responsible for heat-transfer overshoot.


\subsection{Early disturbance amplification}
The streamwise evolution of wall-bounded quantities provides diagnostics for characterising the onset and progression of boundary layer transition. The skin-friction coefficient and Stanton number are defined respectively as
\begin{equation}
    C_f = \cfrac{2\tau_w}{\rho_e u_e^2},
\end{equation}
\begin{equation}
    St = \cfrac{q_w}{\rho_e u_e c_{p,tr} (T_{aw}-T_w)},
\end{equation}
where $T_{aw}$ is the adiabatic wall temperature derived from the respective self-similar solution. For the Stanton number, the translational-rotational temperature $T$ is used for \gls{CPG} and \gls{TNE}, whilst \gls{TEQ} employs a weighted combination of $T$ and $T_v$. The skin-friction coefficient reflects near-wall shear stress and disturbance development, whilst the Stanton number quantifies the nondimensional wall heat flux and thermal boundary-layer response. Both parameters characteristically exhibit overshoot relative to the fully turbulent asymptote during nonlinear breakdown.

For laminar boundary layers, the skin-friction and Stanton number coefficients are derived from self-similarity solutions at each streamwise location to provide a baseline for assessing disturbance growth. Following \citet{White2005}, these are expressed as
\begin{equation}
c_{f,lam} = \sqrt{2} f''(0) \,Re_x^{-1/2}\cfrac{\mu_w \rho_w}{\mu_e \rho_e},
\end{equation}
where $f$ and $\theta$ are the dimensionless velocity and temperature similarity variables. Both depend on the thermal model; formulations are derived from their respective similarity solutions for \gls{CPG}, \gls{TEQ}, and \gls{TNE}.

For turbulent boundary layers, predictions are benchmarked against the compressible Van Driest II correlation \citep{VanDriest1952,Sandhametal2014}:
\begin{equation}
    c_{f,turb} = \frac{1}{\lambda} \left[ \frac{0.242 \left( \arcsin\alpha + \arcsin\beta \right)}{\log_{10}(c_{f,turb}  Re_x) - 0.76 \log_{10}(T_w/T_e) + 0.41} \right],
\end{equation}
where $\lambda = (\gamma - 1) M_e^2 r/2$, $\alpha = (2a^2 - b)/Q$, $\beta = b/Q$, $Q = \sqrt{b^2 + 4a^2}$, $a = \sqrt{\lambda T_e/T_w}$, and $b = (1+\lambda)T_e/T_w - 1$. The Stanton number for both laminar and turbulent counterparts is obtained via the Reynolds analogy:
\begin{equation}
St = \frac{c_f}{2Pr_{tr}^{2/3}}.
\end{equation}

\begin{figure}[t]
    \centering
    \fontsize{9}{10}\selectfont
    \includegraphics[width=0.70\textwidth]{resources/figures/breakdown/result_wallquantities_cf_st.pdf}
    \caption[Streamwise evolution of spatially temporally averaged $C_f$ and $St$ for the Mach~6, $T_w=1800$ K case. Results are shown for CPG, TEQ, and TNE formulations, alongside laminar reference solutions and the Van Driest~II turbulent correlation.]%
    {Streamwise evolution of spatially temporally averaged $C_f$ and $St$ for the Mach~6, $T_w=1800$ K case. Results are shown for \gls{CPG}, \gls{TEQ}, and \gls{TNE} formulations, alongside laminar reference solutions and the Van Driest~II turbulent correlation.}
    \label{figure:thermal_cfst}
\end{figure}

In the present work, laminar and turbulent reference profiles are shown only for the \gls{TNE} case, to avoid excessive clutter in the plots. Comparisons indicate that the laminar and turbulent trends for \gls{CPG} and \gls{TEQ} are very similar to those for \gls{TNE}, although the \gls{TEQ} profiles lie slightly lower. Therefore, the use of \gls{TNE} profiles as representative references is sufficient. 

The skin friction and Stanton number across all cases are shown in \Cref{figure:thermal_cfst}. In the laminar region, all three models closely follow the laminar reference with minor thermal-model differences. A pronounced second-mode resonance peak emerges, with \gls{TEQ} exhibiting a sharp, well-defined peak, \gls{CPG} showing a similar structure delayed by a few displacement thicknesses with slightly higher magnitude, and \gls{TNE} displaying the most attenuated peak. This progressive delay and attenuation reflects different base flows and vibrational relaxation extracting energy from the acoustic modes driving second-mode growth, consistent with~\citep{DiRenzo2021} though not observed in~\citep{Passiatore2021,Williamsetal2025}.

Nonlinear breakdown then accelerates, with both quantities overshooting the turbulent asymptote. \gls{TEQ} exhibits the most dramatic overshoot followed by \gls{CPG} at a downstream location, whilst \gls{TNE} shows considerably subdued overshoot with a more gradual transition to turbulence. The transition hierarchy is clear: \gls{TEQ} transitions earliest, \gls{CPG} transitions downstream, and \gls{TNE} transitions latest, delayed by approximately $150\delta^*$ relative to \gls{TEQ}. The contrasting dynamics reflect both the base-flow effects and the thermal model differences, \gls{TEQ} and \gls{CPG} exhibit sharp transitions with pronounced overshoots followed by rapid relaxation, whilst \gls{TNE} shows a prolonged, gentler approach with suppressed overshoot, demonstrating the stabilising influence of thermal nonequilibrium.

Selected flow quantities and grid parameters at representative stations throughout the transition process are presented in \Cref{table:global_quantities_models}, documenting streamwise evolution across five characteristic phases: near resonance of the two-dimensional second-mode, secondary instability growth, imminent breakdown, end of transition, and the fully turbulent state. Boundary-layer integral parameters including the shape factor $H$ and momentum-thickness Reynolds number $Re_\theta$ characterise velocity profile development, showing progressive steepening as the flow transitions from laminar to turbulent. Thermal quantities including maximum temperature rise $T_{max}/T_e$ quantify thermal excitation effects. In \gls{TNE} flows, the energy balances $B_q = q_w/(\rho_w u_\tau h_w)$ and $B_{q,\mathrm{v}} = q_{w,\mathrm{v}}/(\rho_w u_\tau h_w)$ reveal opposing signs characteristic of energy partitioning between the modes. A clear delay in transition stations is evident across the thermochemical models, with \gls{TNE} consistently requiring greater streamwise distances to progress through successive phases compared to the perfect-gas assumption. The friction Reynolds number $Re_\tau$ and friction Mach number $M_\tau$ increase substantially during breakdown, reflecting intensified near-wall velocity gradients as coherent structures develop and break down. Grid spacings satisfy established \gls{DNS} resolution requirements for capturing near-wall viscous structures and streamwise coherent features, consistent with previous studies~\citep{PoggieEtAl2015}. Similar resolutions were recently achieved using the TENO scheme~\citep{Williamsetal2025}.

\newpage
\begin{sidewaystable}[H]
\centering
% \small
\caption{Streamwise evolution of span-averaged and time-averaged dimensionless flow quantities across the transition process for Mach 6, $T_w = 1800$ K CPG, TEQ and TNE simulations.}
\begin{tabular}{lccc ccc ccc ccc ccc}
\toprule
 & \multicolumn{3}{c}{Near resonance of 2D mode} 
 & \multicolumn{3}{c}{Secondary instability} 
 & \multicolumn{3}{c}{Imminent breakdown} 
 & \multicolumn{3}{c}{End of transition} 
 & \multicolumn{3}{c}{Turbulent} \\
\cmidrule(lr){2-4} \cmidrule(lr){5-7} \cmidrule(lr){8-10} \cmidrule(lr){11-13} \cmidrule(lr){14-16}
 & CPG & TEQ & TNE & CPG & TEQ & TNE & CPG & TEQ & TNE & CPG & TEQ & TNE & CPG & TEQ & TNE \\
\midrule
$x/\delta^*_0$ & 256 & 252 & 292 & 312 & 291 & 361 & 400 & 337 & 436 & 505 & 388 & 505 & 789 & 600 & 600 \\
$Re_x \times 10^{-6}$ & 4.16 & 4.29 & 4.53 & 4.72 & 4.68 & 5.22 & 5.60 & 5.14 & 5.97 & 6.65 & 5.65 & 6.66 & 9.49 & 7.77 & 7.61 \\
$Re_\tau$ & 138 & 145 & 189 & 160 & 151 & 166 & 387 & 175 & 240 & 514 & 254 & 514 & 913 & 687 & 824 \\
$Re^*_\delta \times 10^{-4}$ & 1.61 & 1.56 & 1.68 & 1.72 & 1.63 & 1.80 & 1.87 & 1.71 & 1.93 & 2.04 & 1.79 & 2.04 & 2.43 & 2.10 & 2.18 \\
$Re_\theta$ & 1177 & 1132 & 1211 & 1365 & 1200 & 1333 & 1815 & 1294 & 1537 & 2469 & 1437 & 1809 & 4974 & 2123 & 2323 \\
$H$ & 13.62 & 13.22 & 13.96 & 12.20 & 13.11 & 14.00 & 10.64 & 12.59 & 12.69 & 10.33 & 11.73 & 11.03 & 10.24 & 10.04 & 10.71 \\
$M_{\tau,\mathrm{max}}$ & 0.10 & 0.08 & 0.09 & 0.11 & 0.08 & 0.07 & 0.16 & 0.08 & 0.09 & 0.15 & 0.09 & 0.14 & 0.14 & 0.11 & 0.15 \\
$B_{q,\mathrm{tr}} \times -10^{-2}$ & 7.17 & 2.93 & 5.11 & 6.23 & 2.88 & 3.78 & 8.98 & 3.00 & 4.72 & 8.24 & 3.41 & 6.99 & 7.37 & 4.21 & 7.23 \\
$B_{q,\mathrm{v}} \times 10^{-3}$ & 0.00 & $-5.87$ & 5.58 & 0.00 & $-5.77$ & 5.79 & 0.00 & $-6.01$ & 5.64 & 0.00 & $-6.83$ & 5.69 & 0.00 & $-8.45$ & 5.43 \\
$T_{max}/T_e$ & 3.70 & 3.36 & 3.73 & 3.70 & 3.37 & 3.67 & 3.69 & 3.40 & 3.75 & 3.59 & 3.48 & 3.81 & 3.56 & 3.51 & 3.72 \\
$\delta_{M_\tau}/\delta^{*}_{0}$ & 0.88 & 0.90 & 0.90 & 0.91 & 0.90 & 0.90 & 0.47 & 0.89 & 0.86 & 0.37 & 0.87 & 0.67 & 0.48 & 0.66 & 0.33 \\
\midrule
$\Delta x^+$ & 7.72 & 5.49 & 6.56 & 7.99 & 5.47 & 5.31 & 12.14 & 5.88 & 6.72 & 11.47 & 6.67 & 10.05 & 10.66 & 8.35 & 10.97 \\
$\Delta y_w^+$ & 0.68 & 0.48 & 0.58 & 0.70 & 0.48 & 0.47 & 1.07 & 0.52 & 0.59 & 1.01 & 0.59 & 0.88 & 0.94 & 0.73 & 0.96 \\
$\Delta y^+_{99}$ & 1.54 & 1.52 & 1.97 & 1.74 & 1.58 & 1.72 & 4.00 & 1.81 & 2.46 & 5.22 & 2.59 & 5.20 & 9.15 & 6.89 & 8.27 \\
$\Delta z^+$ & 3.59 & 2.55 & 3.05 & 3.71 & 2.54 & 2.47 & 5.63 & 2.73 & 3.12 & 5.33 & 3.10 & 4.67 & 4.95 & 3.88 & 5.09 \\
$L_z^+$ & 1004 & 714 & 853 & 1038 & 711 & 690 & 1578 & 764 & 873 & 1491 & 867 & 1306 & 1385 & 1085 & 1426 \\
\bottomrule
\end{tabular}
\label{table:global_quantities_models}
\end{sidewaystable}
\newpage

\begin{figure}[H]
    \centering
    \includegraphics[width=22.3cm,angle=-90]{resources/figures/breakdown/result_dns_thermal_stanton_contour.pdf}
    \caption{Time-averaged wall Stanton number ($St$) during boundary layer transition for Mach 6, $T_w = 1800$ K case (a) CPG, (b) TEQ, and (c) TNE models.}
    \label{figure:thermal_spanwise_st}
\end{figure}

The spanwise organisation of heat transfer during transition, as shown in \Cref{figure:thermal_spanwise_st}, reveals noticeable differences across the three thermochemical models. The figure displays contours of Stanton number ($St$) as a function of streamwise position and spanwise coordinate, spanning the transitional region between $x/\delta^*_0 = 200$ and $550$. Streak-like structures emerge as transition progresses, redistributing heat flux into alternating high- and low-intensity bands that mark secondary instability growth and rapid breakdown, consistent with observations in prior transition studies \citep{Franko2013}. The spatial organisation and intensity of these streaks differ markedly across models. \gls{TEQ} develops well-defined streaks early in transition, whilst \gls{CPG} exhibits comparable streak organisation with marginally higher peak intensities during breakdown. Both models maintain organised patterns well into the turbulent region. \gls{TNE} exhibits markedly delayed and attenuated streak development, with high-temperature regions emerging considerably further downstream and remaining less intense than in either \gls{CPG} or \gls{TEQ}. By the end of transition, streak organisation weakens across all models and wall heat flux becomes more uniformly distributed, consistent with fully developed turbulence. This progression from well-organised streaks in \gls{TEQ} and \gls{CPG} to attenuated, delayed streaks in \gls{TNE} reflects the stabilising influence of vibrational nonequilibrium on coherent-structure formation and wall heat transfer.

% \subsection{Transition structures}
% The initial phase of transition is characterised by the amplification of two-dimensional disturbances dominated by the Mack second mode. This behaviour is evident in the modest elevation of skin-friction and Stanton number upstream of breakdown, marking the resonance of the primary instability wavepacket. Linear stability theory predictions presented in Chapter~\ref{chapter:disturbance_growth} demonstrate that the second mode exhibits the largest amplification rates at the frequencies present in this flow. \Cref{figure:lstdns_amplitude} presents linear amplitude growth curves for the representative cases studied herein. The forced frequencies $f = 0.15$ and $f = 0.16$ exhibit peak amplitude growth at distinct streamwise locations, with the observed resonance peaks in the wall-bounded quantities aligning closely with these linear growth regions. This alignment confirms the transition from linear to weakly nonlinear dynamics as energy from the primary modes redistributes into secondary instabilities.

% Instantaneous snapshots reveal the coherent structures driving transition. Firstly, the Q-criterion isosurfaces coloured by total density $\rho / \rho_e$, with pressure contours $p/(\rho_e \ u_e^2 )$ in the background, show how instability modes develop and interact across the three thermal models, as shown in \Cref{figure:transition_isosurfaces}. Initially, two-dimensional Mack modes develop as elongated streamwise-aligned structures. 

% \begin{figure}[H]
%     \centering
%     \includegraphics[width=8.6cm,angle=0]{resources/figures/test/full_isosurfaces.pdf}
%     \caption[Q-criterion isosurfaces coloured by the total density $\rho / \rho_{e}$ during transition for CPG, TEQ, and TNE models.]%
%     {Q-criterion isosurfaces coloured by density during transition for (a) \gls{CPG}, (b) \gls{TEQ}, and (c) \gls{TNE} models. Pressure field shown in background plane. Isosurfaces reveal two-dimensional Mack mode development in the initial domain region and subsequent three-dimensional modal interactions driving breakdown.}
%     \label{figure:transition_isosurfaces}
% \end{figure}


% The resulting resonance rise observed in skin-friction and Stanton-number distributions represents the transition from linear to weakly nonlinear dynamics, transferring energy from the primary mode into secondary instabilities. The secondary structures drive the streak formation \citep{Klebanoff1962} and localised heat-transfer intensification seen in the spanwise distributions, with the strength and organisation of these streaks determining the character of the heat-flux patterns during breakdown.

% The three-dimensional structure of the boundary layer during transition is further illustrated through instantaneous streamwise Q-criterion isosurfaces presented in \Cref{figure:transition_isosurfaces}. The Q-criterion directly reveals the development of low-speed and high-speed streaks characteristic of secondary instability growth across the three thermal models. H-type transition, the most commonly observed phenomenon in hypersonic boundary layers \citep{DeTullio2013}, is driven by the destabilisation of two-dimensional Mack modes \citep{Fedorov2011} and their interaction with oblique second modes \citep{Franko2013}. The \gls{TEQ} case exhibits pronounced, coherent streak formation beginning at comparatively early streamwise locations, with alternating regions of reduced and elevated streamwise velocity apparent at early downstream positions. This vigorous streak development is consistent with rapid secondary instability growth driven by sustained forcing at $f = 0.15$ and $f = 0.16$. The \gls{CPG} case shows similar streak structures but with emergence occurring further downstream, reflecting slightly delayed transition onset. In contrast, the \gls{TNE} case exhibits considerably delayed and attenuated streak development, with velocity fluctuations remaining poorly organised through a substantial portion of the domain. Whilst \gls{TEQ} and \gls{CPG} exhibit enhanced mode interactions leading to direct breakdown to turbulence, the \gls{TNE} case displays characteristic $\Lambda$-vortex structures consistent with observations in adiabatic hypersonic flows \citep{Passiatore2021}, demonstrating the stabilising effect of vibrational nonequilibrium on secondary instability growth and coherent structure formation.




\subsection{Transition structures}

The initial phase of transition is characterised by amplification of two-dimensional disturbances dominated by the Mack second mode. Linear stability theory predictions presented in \Cref{chapter:disturbance_growth} demonstrate that the second mode exhibits the largest amplification rates at the frequencies present in this flow. \Cref{figure:lstdns_amplitude} presents linear amplitude growth curves for the representative cases studied herein. The forced frequencies $f = 0.15$ and $f = 0.16$ exhibit peak amplitude growth at distinct streamwise locations, with the observed resonance peaks in the wall-bounded quantities aligning closely with these linear growth regions. This alignment confirms the transition from linear to weakly nonlinear dynamics as energy from the primary modes redistributes into secondary instabilities.

Instantaneous snapshots reveal the coherent structures driving transition. Q-criterion isosurfaces coloured by total density $\rho / \rho_e$, with pressure contours $p/(\rho_e u_e^2)$ in the background (\Cref{figure:transition_isosurfaces}), show how instability modes develop and interact across the three thermal models. Initially, two-dimensional Mack modes develop as elongated streamwise-aligned structures. These span-aligned structures undergo resonance amplification at the linear growth regions predicted by linear stability theory, transferring energy from the primary mode into secondary instabilities. The secondary structures drive streak formation~\citep{Klebanoff1962} and localised heat-transfer intensification seen in the spanwise distributions, with the strength and organisation of these streaks determining the character of the heat-flux patterns during breakdown.

\begin{figure}[H]
    \centering
    \includegraphics[width=9.6cm,angle=0]{resources/figures/test/full_isosurfaces.pdf}
    \caption{Q-criterion isosurfaces coloured by density during transition for (a) CPG, (b) TEQ, and (c) TNE models. Pressure field shown in background plane. Isosurfaces reveal two-dimensional Mack mode development in the initial domain region and subsequent three-dimensional modal interactions driving breakdown.}
    \label{figure:transition_isosurfaces}
\end{figure}

The three-dimensional boundary-layer structure during transition is illustrated through instantaneous streamwise Q-criterion isosurfaces, which reveal the development of low-speed and high-speed streaks characteristic of secondary instability growth across the three thermal models. H-type transition, the most commonly observed phenomenon in hypersonic boundary layers~\citep{DeTullio2013}, is driven by destabilisation of two-dimensional Mack modes~\citep{Fedorov2011} and their interaction with oblique second modes~\citep{Franko2013}. The \gls{TEQ} case exhibits pronounced, coherent streak formation beginning at comparatively early streamwise locations, with alternating regions of reduced and elevated streamwise velocity apparent at early downstream positions. This vigorous streak development reflects rapid secondary instability growth driven by sustained forcing. The \gls{CPG} case shows similar streak structures but emerging further downstream, reflecting slightly delayed transition onset. In contrast, the \gls{TNE} case exhibits considerably delayed and attenuated streak development, with velocity fluctuations remaining poorly organised through a substantial portion of the domain. Whilst \gls{TEQ} and \gls{CPG} exhibit enhanced mode interactions leading to direct breakdown to turbulence, the \gls{TNE} case displays characteristic $\Lambda$-vortex structures consistent with observations in adiabatic hypersonic flows~\citep{Passiatore2021}, demonstrating the stabilising effect of vibrational nonequilibrium on secondary instability growth and coherent structure formation.

\begin{figure}[t]
    \centering
    \includegraphics[width=18cm,angle=-90]{resources/figures/breakdown/result_contour_yplus1_uvelocity.pdf}
    \caption[Instantaneous streamwise velocity ($u/u_e$) in $x-z$ plane during boundary layer transition for CPG, TEQ, and TNE models.]%
    {Instantaneous streamwise velocity ($u/u_e$) during boundary layer transition for (a) \gls{CPG}, (b) \gls{TEQ}, and (c) \gls{TNE} models.}
    \label{figure:velocity_snapshot}
\end{figure}

Instantaneous streamwise velocity snapshots in the $x$--$z$ plane at $y^+ = 1$ in the turbulent region are shown in \Cref{figure:velocity_snapshot}, revealing streak-like structures in all three cases with Mack mode interaction visible throughout the transitional region. For \gls{CPG} and \gls{TEQ}, modes develop earlier and reach significant amplitude at comparable streamwise locations, driving the observed wall-quantity overshoot~\citep{SchlatturEtAl2008}. The \gls{TNE} case exhibits delayed mode development at more advanced streamwise locations. Despite these spatial delays, vibrational nonequilibrium modulates the spatial progression of transition without fundamentally altering the underlying mechanisms of mode interaction and secondary instability growth. Coherent structures observed here are consistent with incompressible transition dynamics~\citep{SchoppaHussain2002,SchlatturEtAl2008} and with recent analysis of high-amplitude streak dynamics in hypersonic boundary layers~\citep{CaillaurdEtAl2025}, which provide critical insights into the modal interactions governing initial transition stages and nonlinear pathways to turbulence.


% Instantaneous snapshots reveal the coherent structures driving transition. Q-criterion isosurfaces coloured by total density, with pressure contours in the background, show how instability modes develop and interact across the three thermal models, as shown in \Cref{figure:transition_isosurfaces}. Initially, two-dimensional Mack modes develop as elongated streamwise-aligned structures. These grow in amplitude downstream as linear disturbances. The interaction between the two-dimensional Mack mode and three-dimensional secondary modes produces the resonance rise observed in skin-friction and Stanton-number distributions. This resonance represents the transition from linear to weakly nonlinear dynamics, transferring energy from the primary mode into secondary instabilities. The resulting secondary structures drive the streak formation and localised heat-transfer intensification seen in the spanwise distributions, with the strength and organisation of these streaks determining the character of the heat-flux patterns during breakdown. The three-dimensional breakdown of the initially quasi-two-dimensional laminar flow structure is driven by the sustained periodic forcing at $f = 0.15$ and $f = 0.16$, which continues to amplify the primary modes whilst simultaneously exciting oblique secondary waves. The spanwise organisation of wall Stanton number presented in \Cref{figure:thermal_spanwise_st} provides direct evidence of this process: the \gls{TEQ} case develops sharp, well-defined streaks commencing near $x/\delta_0 \approx 300$, whilst the \gls{TNE} case exhibits significantly delayed and attenuated streak development with high-intensity regions appearing substantially further downstream.


% These grow in amplitude downstream as linear disturbances.  The interaction between the two-dimensional Mack mode and three-dimensional secondary modes produces the resonance rise observed in skin-friction and Stanton-number distributions. This resonance represents the transition from linear to weakly nonlinear dynamics, transferring energy from the primary mode into secondary instabilities. The resulting secondary structures drive the streak formation and localised heat-transfer intensification seen in the spanwise distributions, with the strength and organisation of these streaks determining the character of the heat-flux patterns during breakdown.

% The three-dimensional structure of the boundary layer during transition is further illustrated through instantaneous streamwise velocity snapshots in the $xz$ plane at $y^+ = 1$, presented in \Cref{figure:velocity_snapshot}. The velocity field directly reveals the development of low-speed and high-speed streaks that characterise secondary instability growth across the three thermal models. The \gls{TEQ} case exhibits pronounced, coherent streak formation beginning at comparatively early streamwise locations, with alternating regions of reduced and elevated streamwise velocity apparent by approximately $x/\delta_0 \approx 300$. This early and vigorous streak development is consistent with the rapid secondary instability growth driven by sustained forcing at $f = 0.15$ and $f = 0.16$. The \gls{CPG} case shows similar streak structures but with emergence occurring further downstream, reflecting the slightly delayed transition onset compared to \gls{TEQ}. In contrast, the \gls{TNE} case exhibits considerably delayed and attenuated streak development, with clear velocity fluctuations remaining poorly organised through a substantial portion of the domain. The velocity streaks in \gls{TNE} emerge at substantially more advanced streamwise locations and maintain weaker spatial coherence, demonstrating the stabilising effect of vibrational nonequilibrium on secondary instability growth and coherent structure formation.


% The velocity field directly reveals the development of low-speed and high-speed streaks characteristic of secondary instability growth across the three thermal models. H-type transition, the most commonly observed phenomenon in hypersonic boundary layers \citep{DeTullio2013}, is driven by the destabilisation of two-dimensional Mack modes \citep{Fedorov2011} and their interaction with oblique second modes \citep{Franko2013, Passiatore2021}. The \gls{TEQ} case exhibits pronounced, coherent streak formation beginning at comparatively early streamwise locations, with alternating regions of reduced and elevated streamwise velocity apparent at early downstream positions. This vigorous streak development is consistent with rapid secondary instability growth driven by sustained forcing at $f = 0.15$ and $f = 0.16$. The \gls{CPG} case shows similar streak structures but with emergence occurring further downstream, reflecting slightly delayed transition onset. In contrast, the \gls{TNE} case exhibits considerably delayed and attenuated streak development, with velocity fluctuations remaining poorly organised through a substantial portion of the domain. Whilst \gls{TEQ} and \gls{CPG} exhibit enhanced mode interactions leading to direct breakdown to turbulence, the \gls{TNE} case displays characteristic $\Lambda$-vortex structures consistent with observations in adiabatic hypersonic flows \citep{Passiatore2021}, demonstrating the stabilising effect of vibrational nonequilibrium on secondary instability growth and coherent structure formation.

% The three-dimensional structure of the boundary layer during transition is further illustrated through instantaneous streamwise velocity snapshots in the $x-z$ plane at $y^+ = 1$, presented in \Cref{figure:velocity_snapshot}.

% Comparison across the three thermal models reveals important differences in the spatial evolution of these structures. For \gls{CPG} and \gls{TEQ}, Mack modes develop and reach significant amplitude in the earlier part of the domain, consistent with their earlier transition onset. The interaction between two-dimensional and three-dimensional modes occurs at comparable streamwise locations, driving the observed overshoot in wall quantities. In contrast, the \gls{TNE} case exhibits delayed mode development, with Mack modes emerging further downstream and reaching large amplitude at a more advanced streamwise location. Despite these spatial delays, the underlying mechanisms of mode interaction and secondary instability growth remain consistent across all cases, indicating that vibrational nonequilibrium modulates the spatial progression of transition without fundamentally altering the breakdown pathway. 

% The three-dimensional structure of the boundary layer during transition is further illustrated through instantaneous streamwise velocity snapshots in the $x-z$ plane at $y^+ = 1$, presented in \Cref{figure:velocity_snapshot}. Streak-like structures emerge in all three cases, with the effects of Mack mode interaction visible throughout the transitional region. Comparison across the three thermal models reveals important differences in the spatial evolution of these structures. For \gls{CPG} and \gls{TEQ}, Mack modes develop and reach significant amplitude in the earlier part of the domain, consistent with their earlier transition onset. The interaction between two-dimensional and three-dimensional modes occurs at comparable streamwise locations, driving the observed overshoot in wall quantities \citep{SchlatturEtAl2008}. In contrast, the \gls{TNE} case exhibits delayed mode development, with Mack modes emerging further downstream and reaching large amplitude at a more advanced streamwise location. Despite these spatial delays, the underlying mechanisms of mode interaction and secondary instability growth remain consistent across all cases, indicating that vibrational nonequilibrium modulates the spatial progression of transition without fundamentally altering the breakdown pathway. The coherent structures observed herein are consistent with prior studies of transition dynamics \citep{SchoppaHussain2002}. The study by \cite{CaillaurdEtAl2025} provides critical insights into the linear dynamics of hypersonic boundary layers modulated by high-amplitude, steady streaks, mechanisms known to govern the initial stages of transition and dictate subsequent nonlinear pathways to turbulence. Understanding why transition to turbulence is delayed in the \gls{TNE} case requires detailed analysis of the modal interactions and their nonlinear evolution. 



% The differences in modal amplitude and spatial extent reflect the moderated growth rates and stabilising effects of vibrational relaxation discussed in the linear stability analysis, directly linking the instantaneous flow structures to the mean wall-bounded quantities observed in Chapter~\ref{chapter:disturbance_growth}.




% The three-dimensional structure of the boundary layer during transition is further illustrated through instantaneous streamwise velocity snapshots in the $xz$ plane at $y^+ = 1$, presented in \Cref{figure:velocity_snapshot}. The velocity field directly reveals the development of low-speed and high-speed streaks that characterise secondary instability growth across the three thermal models. The \gls{TEQ} case exhibits pronounced, coherent streak formation beginning at comparatively early streamwise locations, with alternating regions of reduced and elevated streamwise velocity apparent by approximately $x/\delta_0 \approx 300$. This early and vigorous streak development is consistent with the rapid secondary instability growth driven by sustained forcing at $f = 0.15$ and $f = 0.16$. The \gls{CPG} case shows similar streak structures but with emergence occurring further downstream, reflecting the slightly delayed transition onset compared to \gls{TEQ}. In contrast, the \gls{TNE} case exhibits considerably delayed and attenuated streak development, with clear velocity fluctuations remaining poorly organised through a substantial portion of the domain. The velocity streaks in \gls{TNE} emerge at substantially more advanced streamwise locations and maintain weaker spatial coherence, demonstrating the stabilising effect of vibrational nonequilibrium on secondary instability growth and coherent structure formation.

% The pronounced differences in transition progression across thermal models reflect fundamental distinctions in how energy redistributes during breakdown. The \gls{TEQ} case undergoes rapid and intense nonlinear breakdown, with the pronounced overshoot in both $C_f$ and $St$ observed approximately $x/\delta_0 \approx 350$--$400$ reflecting swift energy transfer from organised two-dimensional structures to small-scale turbulent features. The \gls{CPG} case exhibits intermediate behaviour with delayed but substantial breakdown dynamics. The \gls{TNE} case demonstrates the most stabilised response, exhibiting delayed and attenuated secondary instability growth with a more gradual transition to turbulence. In the \gls{TNE} case, the vibrational energy mode acts as a thermal energy sink, withdrawing energy from the acoustic modes driving transition. This redistribution into vibrational degrees of freedom moderates the nonlinear breakdown process, resulting in lower transient heating rates and suppressed overshoot intensity. These observations demonstrate that thermal nonequilibrium substantially alters not only the streamwise location of transition but fundamentally modifies the character of the transition process itself, with vibrational relaxation serving as a stabilising mechanism that suppresses the intensity of coherent structure formation and delays the onset of nonlinear breakdown.

% \newpage

% \newpage




% \begin{figure}[H]
%     \centering
%     \includegraphics[width=10.3cm,angle=0]{resources/figures/test/full_isosurfaces.pdf}
%     \caption{Linear stability analysis at $\delta_{*}^{0} = 255$ for (a) CPG, (b) TEQ, and (c) TNE models.}
%     \label{figure:thermal_spanwise_st}
% \end{figure}

% \begin{figure}[H]
%     \centering
%     \includegraphics[width=10.3cm,angle=0]{resources/figures/test/test_plot.png}
%     \caption{Linear stability analysis at $\delta_{*}^{0} = 255$ for (a) CPG, (b) TEQ, and (c) TNE models.}
%     \label{figure:thermal_spanwise_st}
% \end{figure}


% \newpage

% \begin{figure}[H]
%     \centering
%     \includegraphics[width=22.3cm,angle=-90]{resources/figures/breakdown/result_contour_yplus1_uvelocity.pdf}
%     \caption[Instantaneous ($u/u_e$) during boundary layer transition for (a) CPG, (b) TEQ, and (c) TNE models.]%
%     {Instantaneous ($u/u_e$)  during boundary layer transition for (a) \gls{CPG}, (b) \gls{TEQ}, and (c) \gls{TNE} models.}
%     \label{figure:thermal_spanwise_st}
% \end{figure}

% \newpage



% \subsection{Forcing interaction}



% \subsubsection{Boundary layer evolution}

% The development of the boundary layer during transition provides insight into changes in flow structure and momentum redistribution as instability modes grow and break down. Integral quantities such as the displacement thickness $\delta^\ast$ and momentum thickness $\theta^\ast$ characterise the overall profile growth, flow blockage, and momentum deficit. \Cref{figure:dns_thermal_blgrowth} shows the evolution of these quantities for the three thermal models, highlighting how different thermochemical assumptions influence boundary-layer thickening rates and the spatial progression of transition.

% As second-mode resonance develops, growth rates begin to diverge markedly from the laminar regime. In the laminar and early linear instability stages, where vibrational excitation remains modest, the \gls{CPG} and \gls{TNE} predictions exhibit close agreement, reflecting the secondary importance of thermal nonequilibrium effects. However, as transition progresses into the nonlinear breakdown region, a pronounced shift emerges: the \gls{CPG} model transitions towards behaviour more closely aligned with \gls{TEQ}, demonstrating accelerated displacement thickness growth. In contrast, the \gls{TNE} case manifests substantially attenuated thickness development, with thermal relaxation effects increasingly suppressing second-mode growth as vibrational nonequilibrium becomes dynamically significant. By the end of transition ($x/\delta^*_0 > 500$), all three cases undergo rapid boundary layer thickening attributable to coherent structure generation and turbulent cascade processes. The shape factor $H = \delta^* / \theta$ in ~\Cref{table:global_quantities_models} exhibits systematic degradation from approximately $13$--$14$ in the laminar state to below $11$ in fully developed turbulence, reflecting progressive velocity-profile steepening as organised structures develop and dissipate under hypersonic conditions.

% \begin{figure}[t]
%     \centering
%     \includegraphics[width=11cm]{resources/figures/breakdown/result_dns_thermal_boundarylayergrowth.pdf}
%     \caption{Time- and spanwise-averaged boundary layer growth along the streamwise direction for the DNS cases. The evolution of both displacement thickness ($\delta^\ast$), momentum thickness ($\theta^\ast$) and enthalpy thickness ($\delta_{h}$) is shown for the CPG, TEQ, and TNE simulations.}
%     \label{figure:dns_thermal_blgrowth}
% \end{figure}


% This thickening of both displacement and momentum thicknesses accompanies simultaneous reorganisation of the thermal field. The thermal boundary-layer evolution, illustrated through instantaneous temperature distributions in \Cref{figure:thermal_snapshots}, reveals how energy redistribution accompanies momentum-layer development across the three models. In the laminar region (top column), temperature profiles remain smooth and monotonic across all three models, with thermal energy confined to a thin layer adjacent to the wall. As second-mode resonance develops (second column), incipient spanwise temperature modulations emerge, with \gls{TEQ} exhibiting more pronounced thermal striations than \gls{CPG} and \gls{TNE}. During the secondary-instability phase (third column), spanwise temperature variations intensify dramatically, with coherent high-temperature streaks penetrating further into the outer boundary layer. By the nonlinear breakdown stage (last column), complex three-dimensional thermal structures dominate, characterised by intermittent high-temperature filaments and low-temperature pockets reflecting the underlying quasi-streamwise vortex organisation. \gls{TEQ} displays the most vigorous thermal streaking consistent with accelerated transition, whilst \gls{TNE} shows delayed and more attenuated thermal organisation with weaker spanwise modulations. This hierarchy reflects the competing modulation of thermal structures by vibrational nonequilibrium, demonstrating that molecular energy redistribution fundamentally alters both momentum and thermal-structure development during hypersonic boundary-layer transition.

\subsubsection{Boundary layer evolution}

The development of the boundary layer during transition provides insight into changes in flow structure and momentum redistribution as instability modes grow and break down. Integral quantities such as the displacement thickness $\delta^*$ and momentum thickness $\theta^*$ characterise the overall profile growth and momentum deficit. \Cref{figure:dns_thermal_blgrowth} shows the evolution of these quantities for the three thermal models, highlighting how different thermochemical assumptions influence boundary-layer thickening rates and the spatial progression of transition.

As second-mode resonance develops, growth rates diverge markedly from the laminar regime. In the laminar and early linear instability stages, where vibrational excitation remains modest, \gls{CPG} and \gls{TNE} predictions exhibit close agreement, reflecting the secondary importance of thermal nonequilibrium effects. However, as transition progresses into the nonlinear breakdown region, a pronounced shift emerges: the \gls{CPG} model transitions towards behaviour more closely aligned with \gls{TEQ}, demonstrating accelerated displacement thickness growth. In contrast, the \gls{TNE} case manifests substantially attenuated thickness development, with thermal relaxation effects increasingly suppressing second-mode growth as vibrational nonequilibrium becomes dynamically significant. By the end of transition ($x/\delta^*_0 > 500$), all three cases undergo rapid boundary-layer thickening attributable to coherent structure generation and turbulent cascade processes. The shape factor $H = \delta^* / \theta$ in \Cref{table:global_quantities_models} exhibits systematic degradation from approximately $13$--$14$ in the laminar state to below $11$ in fully developed turbulence, reflecting progressive velocity-profile steepening as organised structures develop and dissipate.

\begin{figure}[t]
    \centering
    \includegraphics[width=11cm]{resources/figures/breakdown/result_dns_thermal_boundarylayergrowth.pdf}
    \caption[{Time- and spanwise-averaged boundary-layer growth along the streamwise direction for the DNS cases showing displacement thickness ($\delta^*$), momentum thickness ($\theta^*$), and enthalpy thickness ($\delta_h$) for CPG, TEQ, and TNE simulations.}]{Time- and spanwise-averaged boundary-layer growth along the streamwise direction for the DNS cases showing displacement thickness ($\delta^*$), momentum thickness ($\theta^*$), and enthalpy thickness ($\delta_h$) for \gls{CPG}, \gls{TEQ}, and \gls{TNE} simulations.}
    \label{figure:dns_thermal_blgrowth}
\end{figure}

This thickening of displacement and momentum thicknesses accompanies simultaneous reorganisation of the thermal field. The thermal boundary-layer evolution, illustrated through instantaneous temperature contours in \Cref{figure:thermal_snapshots}, reveals how energy redistribution accompanies momentum-layer development across the three models. From laminar through breakdown stages, temperature contours remain smooth and monotonic in the laminar region across all three models, with thermal energy confined to a thin layer adjacent to the wall. As second-mode resonance develops, incipient spanwise temperature modulations emerge, with \gls{TEQ} exhibiting more pronounced thermal striations than \gls{CPG} and \gls{TNE}. During the secondary-instability phase, spanwise temperature variations intensify dramatically, with coherent high-temperature streaks penetrating further into the outer boundary layer. By the nonlinear breakdown stage, complex three-dimensional thermal structures dominate, characterised by intermittent high-temperature filaments and low-temperature pockets reflecting the underlying quasi-streamwise vortex organisation. \gls{TEQ} displays the most vigorous thermal streaking consistent with accelerated transition, whilst \gls{TNE} shows delayed and more attenuated thermal organisation with weaker spanwise modulations, demonstrating that molecular energy redistribution fundamentally alters both momentum and thermal-structure development during hypersonic boundary-layer transition.

\newpage
\begin{figure}[H]
    \centering
    \includegraphics[width=20.0cm,angle=-90]{resources/figures/breakdown/result_breadkdown_thermal_Tslices.pdf}
    \caption[Instantaneous temperature ($T/T_e$) contours in the $x$--$y$ plane during boundary layer transition for CPG, TEQ, and TNE models.]%
    {Instantaneous temperature ($T/T_e$) contours in the $x$--$y$ plane at a various spanwise locations during nonlinear breakdown and secondary-instability phases for (a) \gls{CPG}, (b) \gls{TEQ}, and (c) \gls{TNE} models. Rows denote progression through the transition region.}
    \label{figure:thermal_snapshots}
\end{figure}
\newpage



\begin{figure}[t]
    \centering
    \includegraphics[width=11.3cm]{resources/figures/breakdown/result_dns_thermal_fft.pdf}
    \caption{Normalised power spectral density and pressure RMS distribution across the three thermochemical models. Panel (a) shows the forced-frequency response and panel (b) shows the spatial evolution of pressure fluctuation intensity during transition.}
    \label{figure:thermal_fft}
\end{figure}

\subsubsection{Frequency content}

FFT analysis was performed on wall pressure fluctuations in the $x$--$z$ plane during statistically steady flow phases across all three models to understand the effects of Mack modes in the transition process. The frequency spectra shown in \Cref{figure:thermal_fft}(a) reveal the response of each thermal model to the dual-frequency excitation. In the laminar region, pressure fluctuations remain minimal with power concentrated at low frequencies. As transition progresses and second-mode resonance develops, distinct spectral peaks emerge at the forced frequencies, with relative amplitudes indicating preferential mode amplification by boundary-layer stability characteristics. Marked spectral broadening is observed compared to the linear predictions of \Cref{chapter:disturbance_growth}, indicating nonlinear energy transfer driving coherent-structure breakdown.

The pressure RMS distribution shown in \Cref{figure:thermal_fft}(b) shows significant differences across the three models. During resonance growth, \gls{CPG} exhibits the highest pressure RMS intensity, whilst \gls{TEQ} and \gls{TNE} display comparatively attenuated fluctuations. However, during secondary instability and breakdown, pressure RMS intensity increases substantially in the \gls{TNE} case, approaching or exceeding that of \gls{CPG}. This inversion of the RMS hierarchy demonstrates the delayed but vigorous nonlinear amplification characteristic of thermally nonequilibrium flows.

% The wall pressure field is analysed through complementary methods to understand the modal mechanisms and spatial organisation during transition: fast Fourier transform (FFT) applied to span-averaged pressure slices to identify dominant frequencies, spatial distribution of pressure RMS to characterise the streamwise evolution of fluctuation intensity, and proper orthogonal decomposition (POD) applied to two-dimensional spanwise-varying pressure data to decompose coherent structures.

% Two Mack second-mode frequencies are excited in the forcing, selected from linear stability theory predictions: $f_1 = 0.15$ and $f_2 = 0.16$ (in nondimensional form). The frequency content of wall pressure fluctuations, shown in \Cref{figure:thermal_fft}, reveals the response of each model to these dual-frequency excitation. In the laminar region, pressure fluctuations remain minimal, with power concentrated at low frequencies. As the domain progresses and the second-mode resonance develops, distinct peaks emerge in the power spectral density at frequencies corresponding to the forced modes. Examining the relative amplitudes of the two forcing frequencies provides insight into which mode is preferentially amplified by the boundary layer stability characteristics. \gls{TEQ} exhibits stronger amplification of both frequencies relative to \gls{CPG}, with the spectral peaks more clearly defined and persisting to higher streamwise locations. The dominance of one frequency over the other varies across models and streamwise locations, reflecting the complex interaction between the modal response and the evolving boundary layer structure during the transition process.




% The spatial distribution of pressure RMS along the streamwise direction characterises how the intensity of wall pressure fluctuations evolves during transition. \Cref{figure:fft_pressure} (lower panel) shows $p_{rms}$ distributions for \gls{CPG} and \gls{TEQ}. In the laminar region ($x/\delta^*_0 < 200$), both models exhibit low and comparable pressure RMS levels, consistent with the minimal disturbance amplitudes in this regime. As the second-mode resonance develops around $x/\delta^*_0 \approx 250$--$350$, pressure RMS begins to rise, with \gls{TEQ} showing elevated levels relative to \gls{CPG}. The peak in pressure RMS occurs during the nonlinear breakdown phase ($x/\delta^*_0 \approx 350$--$450$), where secondary instabilities reach maximum amplitude. \gls{TEQ} sustains higher $p_{rms}$ throughout this region, consistent with its stronger modal amplification observed in the FFT. Following the peak, pressure RMS decays gradually in the turbulent region as energy cascades toward smaller scales and distributes across a broader frequency spectrum. The streamwise location of peak pressure RMS differs between models, occurring earlier in \gls{TEQ} and displaced downstream in \gls{CPG}, directly reflecting the transition delay observed in wall-bounded quantities. These spatial distributions confirm that thermal equilibrium enhances the intensity of pressure fluctuations throughout the transition process, consistent with stronger coherent structure formation.


\begin{figure}[t]
    \centering
    \includegraphics[width=9.3cm]{resources/figures/breakdown/result_dns_thermal_pod_energy.pdf}
    \caption{Eigenvalue spectrum and cumulative energy content of the first 40 POD modes extracted from wall pressure fluctuations, demonstrating the concentration of modal energy in dominant coherent structures across CPG, TEQ, and TNE models.}
    \label{figure:pod_energy_variation}
\end{figure}

\subsubsection{Modal decomposition}

\gls{POD} provides a systematic approach to extract dominant coherent structures from flow fields by decomposing a field into orthogonal spatial modes ranked by energy content. The method decomposes a field $\psi(\mathbf{x},t)$ into a sum of spatial modes $\phi_n(\mathbf{x})$ with corresponding temporal coefficients $a_n(t)$, enabling identification of the most energetically significant structures governing flow dynamics. For incompressible flows, snapshot \gls{POD} has become a standard tool in boundary-layer transition studies, revealing the spatial organisation of large-scale coherent motions~\citep{Sirovich1987}. In compressible hypersonic flows, acoustic fluctuations and pressure-wave propagation complicate the interpretation of standard \gls{POD}. \gls{POD} was performed on wall pressure slices extracted in the $x$--$z$ plane at the wall, recorded over one flow-through time at intervals of 100 iterations. A standard singular value decomposition was employed to compute the spatial modes and their corresponding temporal coefficients, isolating the dominant pressure structures driving transition whilst filtering transient acoustic noise inherent to compressible simulations.

\begin{figure}[t]
    \centering
    \includegraphics[width=18.0cm,angle=-90]{resources/figures/breakdown/result_dns_thermal_contour_pod_mode1.pdf}
    \caption{Spatial structure of dominant POD mode extracted from wall pressure fluctuations during boundary-layer transition for (a) CPG, (b) TEQ, and (c) TNE models.}
    \label{figure:contour_podmode1}
\end{figure}

Spectral Proper Orthogonal Decomposition (SPOD) extends this framework by decomposing the flow into frequency-dependent modes, isolating coherent structures at specific frequencies and separating them from broadband acoustic noise \citep{TowneEtAl2018}. However, SPOD requires substantial computational resources and extensive data archives, making snapshot \gls{POD} more practical for \gls{DNS} studies with large datasets. Here, snapshot \gls{POD} is applied to wall pressure fluctuations in the $x$--$z$ plane to characterise the spatial structure and energy distribution of coherent pressure patterns during transition across the three thermochemical models.

% \begin{figure}[t]
%     \centering
%     \includegraphics[width=18.0cm,angle=-90]{resources/figures/breakdown/result_dns_thermal_contour_pod_mode30.pdf}
%     \caption{Spatial structure of the POD mode 30 extracted from wall pressure fluctuations during boundary-layer transition for (a) CPG, (b) TEQ, and (c) TNE models, revealing differences in coherent-structure organisation.}
%     \label{figure:contour_podmode30}
% \end{figure}

% The spatial structure of the leading \gls{POD} mode  seen in \Cref{figure:contour_podmode1}, reveals that coherent structures exhibit similar organisational characteristics across all three thermal models, with spanwise striations present in \gls{CPG}, \gls{TEQ}, and \gls{TNE}. However, the transition zone in \gls{TNE} is notably more spatially compressed in the streamwise direction, with transition-related modal content concentrated over a narrower region compared to the more extended transition zones in \gls{CPG} and \gls{TEQ}. Examination of higher-order modes (\Cref{figure:contour_podmode30}) reveals distinct differences in modal energy distribution. \gls{TEQ} exhibits the highest energetic content in mode 30, whilst \gls{TNE} and \gls{CPG} retain progressively less energy in higher modes. When considered alongside skin friction evolution (\Cref{figure:mach_cf}), which shows \gls{TEQ} and \gls{CPG} exhibiting sharp transitions with steep gradients whilst \gls{TNE} transitions more gradually over an extended streamwise range despite delayed onset, a complex picture emerges. The spatial compression of the forced Mack mode in \gls{TNE} does not correspond to abrupt transition. Rather, \gls{TNE}'s concentrated modal content is followed by prolonged energy transfer to higher modes and gradual skin friction rise. In contrast, \gls{TEQ} displays both an extended primary modal structure and the highest energy in higher modes, yet exhibits the sharpest transition in skin friction. These observations indicate that the relationship between primary instability spatial extent, modal energy cascade, and wall transition characteristics is fundamentally altered by thermal nonequilibrium, reflecting qualitatively distinct transition pathways across the three models. The sharp transition in \gls{TEQ} may reflect rapid nonlinear saturation driven by accumulated linear instability growth coupled with wall-imposed forcing, whilst the gradual transition in \gls{TNE} indicates that vibrational nonequilibrium modifies nonlinear energy-transfer mechanisms. Further investigation of N-factor evolution and wall-forcing coupling across thermal models would clarify these distinctions.

The spatial structure of the leading \gls{POD} mode seen in \Cref{figure:contour_podmode1} reveals that coherent structures exhibit similar organisational characteristics across all three thermal models, with spanwise striations present in \gls{CPG}, \gls{TEQ}, and \gls{TNE}. However, the transition zone in \gls{TNE} is notably more spatially compressed in the streamwise direction, with transition-related modal content concentrated over a narrower region compared to the more extended transition zones in \gls{CPG} and \gls{TEQ}. Examination of higher-order modes in \Cref{figure:contour_podmode30} reveals distinct differences in modal energy distribution. \gls{TEQ} exhibits the highest energetic content in mode 30, whilst \gls{TNE} and \gls{CPG} retain progressively less energy in higher modes. When considered alongside skin friction evolution in \Cref{figure:mach_cf}, which shows \gls{TEQ} and \gls{CPG} exhibiting sharp transitions with steep gradients whilst \gls{TNE} transitions more gradually over an extended streamwise range despite delayed onset, a complex picture emerges.

The spatial compression of the forced Mack mode in \gls{TNE} does not correspond to abrupt transition. Rather, \gls{TNE}'s concentrated modal content is followed by prolonged energy transfer to higher modes and gradual skin friction rise. In contrast, \gls{TEQ} displays both an extended primary modal structure and the highest energy in higher modes, yet exhibits the sharpest transition in skin friction. These observations indicate that the relationship between primary instability spatial extent, modal energy cascade, and wall transition characteristics is fundamentally altered by thermal nonequilibrium, reflecting qualitatively distinct transition pathways across the three models. The sharp transition in \gls{TEQ} may reflect rapid nonlinear saturation driven by accumulated linear instability growth coupled with wall-imposed forcing, whilst the gradual transition in \gls{TNE} indicates that vibrational nonequilibrium modifies nonlinear energy-transfer mechanisms.



% The compressed streamwise extent of mode-1 content in \gls{TNE}, coupled with elevated energy in higher modes, suggests that the forced Mack mode exhibits accelerated nonlinear saturation in the presence of vibrational nonequilibrium. Whilst \gls{CPG} and \gls{TEQ} show gradual modal growth and transition over an extended streamwise region, \gls{TNE} exhibits more rapid cascade to broadband turbulence, indicating that thermal nonequilibrium modifies the modal energy-transfer mechanism during instability breakdown.

% The energy spectrum shown in \Cref{figure:pod_energy_variation} demonstrates that \gls{TNE} concentrates substantially more energy in the dominant POD mode (approximately $0.18$) compared to \gls{CPG} and \gls{TEQ} (approximately $0.12$), indicating a more energetically dominant coherent structure. The cumulative energy distribution shows that \gls{TNE} achieves approximately $60\%$ of total energy with the first few modes, whilst \gls{CPG} and \gls{TEQ} require a more distributed modal contribution to reach equivalent energy levels. This disparity suggests that transition in \gls{TNE} is governed by a single leading coherent mechanism, whereas \gls{CPG} and \gls{TEQ} exhibit more complex, distributed modal interactions. The spatial structure of the dominant mode, presented in \Cref{figure:contour_podmode1}, further illustrates these differences. \gls{CPG} and \gls{TEQ} display organised spanwise striations reflecting continuous streak systems \citep{SchoppaHussain2002}, whilst \gls{TNE} exhibits localised, concentrated pressure perturbations with less continuous organisation. Higher-order modes beyond mode $25$ begin to capture broadband turbulent structures, consistent with the transition to fully developed turbulence. Although not presented for brevity, Fourier analysis of the POD temporal coefficients recovers the fundamental forced frequencies, confirming that the dominant modes directly capture the modal dynamics induced by the imposed excitation. These observations indicate that despite vibrational nonequilibrium delaying transition onset in \gls{TNE}, the resulting coherent structures, when they emerge, are characterised by greater spatial localisation and energetic dominance, a signature of fundamentally different nonlinear dynamics governing the breakdown pathway in \gls{TNE} flow. The delayed transition in \gls{TNE} relative to \gls{TEQ} and \gls{CPG} is thus not merely a shift in streamwise location, but reflects a qualitatively distinct mechanism of coherent-structure formation driven by the stabilising effects of vibrational nonequilibrium.

% The energy spectrum shown in \Cref{figure:pod_energy_variation} demonstrates that \gls{TNE} concentrates substantially more energy in the dominant POD mode (approximately $0.18$) compared to \gls{CPG} and \gls{TEQ} (approximately $0.12$), indicating a more energetically dominant coherent structure. The cumulative energy distribution shows that \gls{TNE} achieves approximately $60\%$ of total energy with the first few modes, whilst \gls{CPG} and \gls{TEQ} require a more distributed modal contribution to reach equivalent energy levels. This disparity suggests that transition in \gls{TNE} is governed by a single dominant coherent mechanism, whereas \gls{CPG} and \gls{TEQ} exhibit more complex, distributed modal interactions. The spatial structure of the dominant mode, presented in \Cref{figure:contour_podmode1}, further illustrates these differences. \gls{CPG} and \gls{TEQ} display organised spanwise striations reflecting continuous streak systems \citep{SchoppaHussain2002}, whilst \gls{TNE} exhibits localised, concentrated pressure perturbations with less continuous organisation. Higher-order modes beyond mode $25$ begin to capture broadband turbulent structures, consistent with the transition to fully developed turbulence~\citep{}. These observations indicate that despite vibrational nonequilibrium delaying transition onset in \gls{TNE}, the resulting coherent structures, when they emerge, are characterised by greater spatial localisation and energetic dominance, a signature of fundamentally different nonlinear dynamics governing the breakdown pathway in \gls{TNE} flow. The delayed transition in \gls{TNE} relative to \gls{TEQ} and \gls{CPG} is thus not merely a shift in streamwise location, but reflects a qualitatively distinct mechanism of coherent-structure formation driven by the stabilising effects of vibrational nonequilibrium.


% To understand the spanwise variation of pressure fluctuations and identify the dominant coherent structures, proper orthogonal decomposition is applied to wall pressure data on a two-dimensional x-z slice (streamwise and spanwise directions) at fixed times. The energy distribution across the first 40 POD modes, shown in \Cref{figure:pod_energy}, reveals significant differences between \gls{CPG} and \gls{TEQ}. The first POD mode contains approximately 11.5\% of the total kinetic energy in \gls{TEQ} compared to 11\% in \gls{CPG}, indicating that the dominant coherent structure is comparable in both cases but carries slightly higher energy in the thermally equilibrated case. However, the cumulative energy distribution diverges beyond the first few modes: \gls{TEQ} reaches 50\% cumulative energy by mode 8--9, whilst \gls{CPG} requires modes up to approximately 15--18 to achieve the same threshold. This difference indicates that \gls{TEQ} organises pressure fluctuations into fewer, more energetic modes, reflecting the strong coherent streak structures observed in the spanwise heat-transfer distributions. In contrast, \gls{CPG} distributes energy across a broader range of modes, suggesting weaker coherent organisation and more fragmented spatial structure during transition. By mode 40, the cumulative energy reaches approximately 54\% in \gls{CPG} and 83\% in \gls{TEQ}, demonstrating that \gls{TEQ} concentrates more of the total pressure-field energy into large-scale modal structures, whilst \gls{CPG} retains significant energy in unresolved small-scale features.

% Proper Orthogonal Decomposition of wall pressure fluctuations reveals fundamental differences in the modal energy distribution and spatial coherence across the three thermochemical models. The energy spectrum shown in \Cref{figure:pod_energy_variation} demonstrates that \gls{TNE} concentrates substantially more energy in the dominant POD mode (approximately $0.18$) compared to \gls{CPG} and \gls{TEQ} (approximately $0.12$), indicating a more energetically dominant coherent structure \citep{Lumley1967, Sirovich1987}. The cumulative energy distribution shows that \gls{TNE} achieves approximately $60\%$ of total energy with the first few modes, whilst \gls{CPG} and \gls{TEQ} require a more distributed modal contribution to reach equivalent energy levels. This disparity suggests that transition in \gls{TNE} is governed by a single dominant coherent mechanism, whereas \gls{CPG} and \gls{TEQ} exhibit more complex, distributed modal interactions. The spatial structure of the dominant mode, presented in \Cref{figure:contour_podmode1}, further illustrates these differences. \gls{CPG} and \gls{TEQ} display organised spanwise striations reflecting continuous streak systems \citep{SchoppaHussain2002}, whilst \gls{TNE} exhibits localised, concentrated pressure perturbations with less continuous organisation. These observations indicate that despite vibrational nonequilibrium delaying transition onset in \gls{TNE}, the resulting coherent structures, when they emerge, are characterised by greater spatial localisation and energetic dominance—a signature of fundamentally different nonlinear dynamics governing the breakdown pathway in thermochemically nonequilibrium flows.

% The Proper Orthogonal Decomposition (POD) analysis decomposes the pressure fluctuations into a set of spatial modes that capture the dominant coherent structures in the flow. The methodology begins by assembling the pressure field data from multiple timesteps into a spatiotemporal matrix of dimensions $(n_x \cdot n_z) \times n_t$, where $n_x$ and $n_z$ are the streamwise and spanwise grid points, and $n_t$ is the number of temporal snapshots. The data is mean-centred by subtracting the temporal average from each spatial location to isolate fluctuations. The Snapshot POD approach employs Singular Value Decomposition (SVD) on this centred matrix, which yields left singular vectors (spatial modes), right singular vectors (temporal coefficients), and singular values. The energy contained in each mode is quantified as $E_i = \lambda_i^2 / \sum \lambda_j^2$, where $\lambda_i$ represents the $i$-th singular value, allowing identification of the modes that capture the most energetic structures. The cumulative energy $CE_n = \sum_{i=1}^{n} E_i$ provides insight into the number of modes required to represent a given fraction of the total energy in the flow field. This decomposition facilitates direct comparison between the CPG and TEQ thermochemical models by examining whether the dominant modes and energy distribution differ significantly between the two approaches, thereby revealing differences in the coherent structure organisation across thermochemical assumptions.




% The structure of POD mode 1, the most energetic single mode, represents the dominant coherent structure in the flow. Spatial analysis of mode 1 reveals quasi-periodic spanwise patterns that correspond directly to the streak structures identified in instantaneous and time-averaged Stanton-number distributions. The mode exhibits strong amplitude in the breakdown region ($x/\delta^{*}_{0} \approx 350$--$450$) and persists downstream into the turbulent region, reflecting the sustained spanwise organisation of heat transfer during transition. The spanwise wavelength captured by mode 1 is approximately $\lambda_z \approx 100$--$120\delta^{*}_{0}$, consistent with the subharmonic wavelengths introduced in the initial forcing and their subsequent nonlinear evolution. The mode 1 amplitude in \gls{TEQ} exceeds that in \gls{CPG}, indicating stronger coherent streak formation driven by the enhanced secondary instability growth in the thermally equilibrated case. In contrast, higher-order modes (e.g., mode 34) capture the turbulent fluctuations that develop in the fully turbulent region and the small-scale pressure variations associated with vortical structures. Mode 34 exhibits rapid spatial variation and lacks the organised periodic structure of mode 1, instead showing the signature of fine-scale turbulent eddies. The energies of mode 34 are comparable between \gls{CPG} and \gls{TEQ}, suggesting that once fully turbulent, the cascade of energy toward small scales proceeds at similar rates regardless of thermal modelling.




% The progression of POD mode energies from mode 1 (large-scale streaks) through the intermediate modes to mode 34 (small-scale turbulence) illustrates the energy cascade mechanism during transition. The steeper energy decay observed in \gls{TEQ} (modes 1--10) compared to \gls{CPG} reflects the more coherent structure and concentration of energy in the large-scale modes. The flatter decay in \gls{CPG} (modes 1--20) indicates that energy is more uniformly distributed across scales, consistent with a less organised transition process. These findings confirm that thermal equilibrium enhances coherent structure formation, concentrating pressure fluctuations into energetic streak modes, whilst perfect-gas dynamics support a more distributed modal response. The POD analysis therefore provides a complementary perspective to the instantaneous isosurface observations and time-averaged Stanton-number distributions, directly linking the spectral properties of wall pressure to the spatial organisation and intensity of coherent structures during breakdown.




% \subsection{Boundary layer evolution}

% This section is where I need to include things about Boundary layer profiles and temperature profiles to show how the steep temperatures are effects the total features, maybe quantify here that the thermal nonequilbrium is pushed? or should I do this in the ILES chapter? 



% The development of the boundary layer during transition provides insight into the changes in flow structure and momentum redistribution as instability modes grow and breakdown occurs. Integral parameters such as displacement thickness and momentum thickness characterise the overall profile development and the degree of flow blockage and momentum deficit. Examining these quantities across the three thermal models reveals how thermal modelling assumptions affect the rate of boundary layer thickening and the spatial distribution of transition stages.

% The streamwise evolution of displacement thickness $\delta^*$ and momentum thickness $\theta$ is shown in \Cref{figure:blgrowth}. In the laminar region ($x/\delta^*_0 < 200$), all three models exhibit similar growth rates, with $\delta^*$ and $\theta$ increasing gradually in accordance with Blasius-type scaling expected for a developing laminar boundary layer. The agreement across models in this region reflects the minimal influence of thermal effects when disturbance amplitudes remain small. As the second-mode resonance develops around $x/\delta^*_0 \approx 250$--$350$, the growth rates begin to diverge. \gls{TEQ} exhibits accelerated thickening consistent with its early transition onset, whilst \gls{CPG} shows intermediate growth. The \gls{TNE} case displays the most gradual thickening in this region, reflecting the delayed breakdown and attenuated instability growth rates predicted by linear studies. By the end of transition ($x/\delta^*_0 > 500$), all three models show accelerated growth as the flow becomes increasingly turbulent. The final displacement and momentum thicknesses converge toward comparable values, though the streamwise locations at which saturation is achieved differ across models, directly reflecting the transition delays observed in wall-bounded quantities. The shape factor $H = \delta^* / \theta$ provides additional insight into profile evolution, transitioning from approximately $13$--$14$ in the laminar state to below $11$ in fully developed turbulence, indicating progressive steepening of the velocity profile as coherent structures develop and breakdown proceeds.

\subsubsection{Disturbance growth}

The evolution of disturbance kinetic energy across the transition process was examined to understand flow behaviour. Disturbance kinetic energy is defined as
\begin{equation}\label{equation:tke}
e_{dke} = \frac{1}{2}\left(\widetilde{u}'^2 + \widetilde{v}'^2 + \widetilde{w}'^2\right),
\end{equation}
where $\widetilde{u}'$, $\widetilde{v}'$, and $\widetilde{w}'$ are the root-mean-square velocity fluctuations obtained from Favre-averaged variables in the streamwise, wall-normal, and spanwise directions, respectively. This quantity directly quantifies the energy contained in velocity perturbations and reveals how instability modes amplify and redistribute energy across spatial regions during transition. \Cref{figure:eke_budget} presents contours of disturbance kinetic energy in the $x$--$y$ plane with time and spanwise averaging, for the Mach 6, $T_w = 1800~\text{K}$ case across \gls{CPG}, \gls{TEQ}, and \gls{TNE} formulations. The aspect ratio is intentionally non-unity to clearly visualise the streamwise evolution of disturbance energy as it concentrates near the wall during linear growth, intensifies during secondary instability, and eventually extends throughout the boundary layer as turbulence develops. 

\begin{figure}[t]
    \centering
    \fontsize{9}{10}\selectfont
    \includegraphics[width=0.80\textwidth]{resources/figures/breakdown/result_dns_thermal_tke_m6tw18.pdf}
    \caption[Time- and spanwise-averaged disturbance kinetic energy contours in the $x$--$y$ plane for Mach 6, $T_w = 1800~\text{K}$ comparing CPG, TEQ, and TNE formulations.]{Time- and spanwise-averaged disturbance kinetic energy contours in the $x$--$y$ plane for Mach 6, $T_w = 1800~\text{K}$ comparing \gls{CPG}, \gls{TEQ}, and \gls{TNE} formulations.}
    \label{figure:eke_budget}
\end{figure}

The disturbance kinetic energy contours presented in \Cref{figure:eke_budget} provide direct quantification of energy evolution across the transition process. For \gls{CPG} and \gls{TEQ}, the spatial extent of energy is distributed across the wall-normal direction early in transition, with distinct peaks occurring at comparable streamwise locations and reflecting rapid transition onset. In contrast, \gls{TNE} exhibits more confined, intense energy concentrations near the wall initially, which extend progressively into the boundary layer at downstream locations. This spatial confinement reflects the delayed amplification of disturbances, yet the ultimately higher energy magnitude in the \gls{TNE} case demonstrates the vigorous nonlinear growth following secondary instability emergence. The wall-normal penetration of disturbance energy differs markedly in the fully developed region, with \gls{TNE} energy extending higher relative to boundary-layer thickness, indicating fundamentally different turbulent structure organisation modulated by vibrational nonequilibrium. The earlier saturation of energy growth in \gls{CPG} and \gls{TEQ} contrasts with the sustained energy rise in \gls{TNE}, reflecting the prolonged transition pathway in thermally nonequilibrium flows. These observations unify the transition behaviour revealed by skin friction, spectral content, modal decomposition, and mean-flow statistics, demonstrating that thermal nonequilibrium systematically delays transition whilst preserving the fundamental instability mechanisms.

\subsection{Turbulent structures}

Once the boundary layer transitions through nonlinear breakdown and secondary-instability phases, thermal structures progressively dissipate into small-scale turbulent fluctuations whilst mean-flow properties asymptotically approach turbulent equilibrium values. In this fully turbulent regime, the three thermochemical models exhibit markedly different turbulence characteristics despite comparable integral boundary-layer growth. To isolate and quantify model-dependent effects on turbulence structure, velocity statistics and Reynolds stress distributions are examined. Statistical averaging was performed over $1.5$ flow-through times following establishment of fully turbulent conditions, corresponding to approximately 50 eddy turnover times~\citep{Pirozzoli2004}, sufficient to capture converged statistics whilst resolving the high-frequency acoustic modes characteristic of hypersonic boundary layers. Representative stations within the turbulent region are analysed for each model, with streamwise locations specified in \Cref{table:global_quantities_models}.

% \subsubsection{Velocity and temperature statistics}
% A detailed assessment of velocity statistics is essential for characterising the turbulent boundary layer. Mean velocity profiles provide direct information on the redistribution of momentum by turbulence and form the basis for testing scaling laws and transformation approaches. Higher-order statistics, such as Reynolds stresses, quantify the intensity and anisotropy of turbulent fluctuations and reveal how energy is transferred amongst flow components. Together, these measures offer a comprehensive view of the turbulent structure and allow differences between thermal models to be identified and interpreted.

% In compressible boundary layers, direct comparison with incompressible velocity profiles requires a transformation that accounts for density and viscosity variations near the wall. The \cite{VanDriest1952} transformation extends the original van Driest scaling to include viscous effects, offering improved collapse of velocity profiles in the inner region. This provides a common reference for assessing compressible turbulent boundary layers and comparing different thermal models on a consistent basis. The mean velocity and wall-normal coordinate $y$ are scaled using the friction velocity $u_{\tau}$ and the viscous length scale $\nu_w/u_{\tau}$. The transformed velocity is defined as
% \begin{equation}
%     u^{+} = \int_{0}^{\bar{u}} \sqrt{\frac{\bar{\rho}}{\rho_w}} \, \mathrm{d}u ,
% \end{equation}
% where $\bar{\rho}$ is the time-averaged local mean density and $\rho_w$ is the density at the wall.

% The resulting mean velocity profile should collapse onto the universal law of the wall, matching the viscous sublayer ($u^{+}=y^{+}$ for $y^{+}<10$) and the logarithmic region at larger wall-normal distances:
% \begin{equation}
%     u^{+} = \frac{1}{\kappa}\ln(y^{+}) + C, 
%     \qquad \kappa = 0.41, \quad C = 5.2 .
% \end{equation}





% Recent work has introduced improved velocity transformations for compressible profiles, particularly under intrinsic compressibility \cite{HasanEtal2023}. These approaches build on semilocal scaling concepts \cite{TrettelLarsson2016,GriffinEtAl2021}, which emphasise the role of variable density and viscosity in defining an effective viscous length scale. \cite{KumarLarsson2022} looking into this adn then finally, ~\cite{HasanEtal2023} explored intrinsic compressiblity effect on the transformation, while  ~\cite{HasanEtAl2024} provided a more suitable transfomration for diverse wall conditions and MAch numbers.  

\subsubsection{Velocity and temperature statistics}
A detailed assessment of velocity statistics is essential for characterising the turbulent boundary layer. Mean velocity profiles provide direct information on the redistribution of momentum by turbulence and form the basis for testing scaling laws and transformation approaches. Higher-order statistics, such as Reynolds stresses, quantify the intensity and anisotropy of turbulent fluctuations and reveal how energy is transferred amongst flow components. Together, these measures offer a comprehensive view of the turbulent structure and allow differences between thermal models to be identified and interpreted.

In compressible boundary layers, direct comparison with incompressible velocity profiles requires a transformation that accounts for density and viscosity variations near the wall. The van Driest transformation \citep{VanDriest1952} extends the original scaling to include density effects, offering improved collapse of velocity profiles in the inner region. This provides a common reference for assessing compressible turbulent boundary layers and comparing different thermal models on a consistent basis. The mean velocity and wall-normal coordinate $y$ are scaled using the friction velocity $u_{\tau}$ and the viscous length scale $\nu_w/u_{\tau}$. The transformed velocity is defined as
\begin{equation}
    u^{+} = \int_{0}^{\bar{u}} \sqrt{\frac{\bar{\rho}}{\rho_w}} \, \mathrm{d}u ,
\end{equation}
where $\bar{\rho}$ is the time-averaged local mean density and $\rho_w$ is the density at the wall.

The resulting mean velocity profile should collapse onto the universal law of the wall, matching the viscous sublayer ($u^{+}=y^{+}$ for $y^{+}<10$) and the logarithmic region at larger wall-normal distances:
\begin{equation}
u^{+} = \frac{1}{\kappa}\ln(y^{+}) + C, 
\qquad \kappa = 0.41, \quad C = 5.2.
\end{equation}

Recent developments have introduced improved velocity transformations for compressible profiles. Semilocal scaling concepts~\citep{TrettelLarsson2016,GriffinEtAl2021,KumarLarsson2022} emphasise the role of variable density and viscosity in defining an effective viscous length scale. Investigations of van Driest transformations accounting for intrinsic compressibility~\citep{HasanEtal2023,HasanEtal2024} have revealed that this effect is a characteristic feature of thermochemical nonequilibrium turbulence~\citep{Williamsetal2025}, whilst cold-wall conditions suppress it~\citep{Zhangetal2018}. The present work employs the transformations suggested by~\citep{HasanEtal2024}, which provide enhanced accuracy across diverse wall conditions and Mach numbers.

Under the assumption of universality of the turbulent shear stress in the near-wall region, the semilocal wall coordinate is related to the wall-normal coordinate through
\begin{equation}
y^{*} = \frac{\delta_{v}}{\delta^{*}_{v}}\, y^{+},
\end{equation}



% To define the transformation, a consistent relationship between $Y^{+}$ and $y^{+}$ is required. Under the assumption of universality of the turbulent shear stress in the near-wall region, the relation
% \begin{equation}
%     Y^{+} = \frac{\delta_{v}}{\delta^{*}_{v}}\, y^{+} = y^{*}
% \end{equation}
% is introduced. In practice, this universality is not guaranteed once compressibility modifies the near-wall stress distribution, so the equivalence $Y^{+}=y^{*}$ requires reassessment. The interpretation
% \begin{equation}
%     Y^{+} = \frac{\nu}{\delta^{*}_{v}}
% \end{equation}
% identifies $\delta^{*}_{v}$ as the relevant viscous length scale for small-scale motions in compressible flows and forms the basis of the HLPP formulation.

% Making use of the coordinate mapping $Y^{+}=y^{*}$, together with
% \begin{equation}
%     \frac{dy^{*}}{dy} = \left(1-y^{*}\frac{d\delta^{*}_{v}}{dy}\right)\frac{\delta_{v}}{\delta^{*}_{v}},
%     \qquad
%     \frac{u_{\tau}}{u^{*}_{\tau}}=\sqrt{\frac{\bar{\rho}}{\rho_{w}}},
% \end{equation}
% the transformation kernel becomes
% \begin{equation}
% \frac{d\overline{U}^{+}}{dy^{+}} = 
% \left(\frac{1+\mu^{c}_{t}/\bar{\mu}}{1+\mu^{i}_{t}/\mu_{w}}\right)^{1/3}
% \left(1-\frac{y}{\delta^{*}_{v}}\frac{d\delta^{*}_{v}}{dy}\right)^{1/2}
% \sqrt{\frac{\bar{\rho}}{\rho_{w}}},
% \label{eq:TLkernel}
% \end{equation}
% where the three multiplicative terms represent adjustments to (i) friction velocity scaling, (ii) the viscous length scale, and (iii) density variations in the inner layer.

% To close the formulation, explicit models for the eddy viscosities $\mu^{i}_{t}$ and $\mu^{c}_{t}$ are introduced. The Johnson–King model with Van Driest-type damping is adopted, giving
% \begin{equation}
%     \mu^{i}_{t}=\sqrt{\tau_{w}\rho_{w}}\,\kappa y D^{i},
%     \qquad
%     D^{i}=\left[1-\exp\left(-\frac{y^{+}}{A^{+}}\right)\right]^{2},
% \end{equation}
% with $\kappa=0.41$ and $A^{+}=17$ yielding a log-law intercept close to 5.2. For compressible flows, the corresponding form is
% \begin{equation}
%     \mu^{c}_{t}=\sqrt{\tau_{w}\rho_{w}}\,\kappa y D^{c},
% \end{equation}
% where $D^{c}$ is a damping function based on the semilocal coordinate $y^{*}$.

% Integrating the kernel gives the transformed velocity profile,
% \begin{equation}
% u^{+}_{HLPP}(y^{*}) = \int_{0}^{y^{*}}
% \left(\frac{1+\mu^{c}_{t}/\bar{\mu}}{1+\mu^{i}_{t}/\mu_{w}}\right)^{1/3}
% \left(1-\frac{y}{\delta^{*}_{v}}\frac{d\delta^{*}_{v}}{dy}\right)^{1/2}
% \sqrt{\frac{\bar{\rho}}{\rho_{w}}}\, dy^{*},
% \label{eq:hlpp}
% \end{equation}
% where $u^{+}_{HLPP}$ denotes the transformed velocity in wall units. This representation allows mean velocity profiles in compressible boundary layers to be compared consistently across thermal models and remains accurate even under cold-wall conditions.


\begin{figure}[t]
    \centering
    \fontsize{9}{10}\selectfont
    \includegraphics[width=0.60\textwidth]{resources/figures/breakdown/result_dns_velocity_profiles.pdf}
    \caption[Velocity and temperature statistics in transformed wall units for Mach 6, $T_w=1800$ K.]{Profiles of velocity and temperature statistics plotted against the semilocal wall coordinate $y^{*}$, obtained from the van Driest transformation, for Mach 6, $T_w=1800$ K comparing \gls{CPG}, \gls{TEQ}, and \gls{TNE} formulations.}
    \label{figure:velocity_statistics}
\end{figure}

The velocity profiles collapse well onto the expected scaling laws when expressed in terms of the semilocal wall coordinate $y^{*}$, following the linear relation in the viscous sublayer and approaching the logarithmic law in the overlap region as shown in \Cref{figure:velocity_statistics}. This confirms that the van Driest transformation provides a consistent framework across the different thermal models. The turbulent Mach number $M_t$ exhibits clear thermal sensitivity, with \gls{TEQ} producing the highest peak, followed by \gls{TNE}, and \gls{CPG} remaining lowest. This ranking mirrors observations in adiabatic flows~\citep{Duan2011}, where increasing freestream Mach number produces similar shifts in $M_t$ profiles. These results demonstrate that vibrational nonequilibrium fundamentally alters turbulent structure, with vibrational excitation effects rivalling those of bulk flow acceleration. Despite elevated disturbance kinetic energy in \gls{TNE}, the turbulent Mach number remains intermediate, indicating that vibrational relaxation redistributes energy towards non-acoustic fluctuation modes rather than enhancing acoustic velocity fluctuations comparable to \gls{TEQ}.


% \cite{FievetEtAl2019} explored the coupling of vibrational nonequilibrium and turbulent mixing, observing that  vibrational nonequlibrium is enhanced by 

% The temperature profiles, however, do not show full equilibration. Whereas linear stability theory indicated that in the downstream region the \gls{TNE} mean temperature tends to align more closely with \gls{TEQ} than with \gls{CPG}, the present turbulent results reveal the opposite behaviour: turbulence slows down the equilibration process, and the \gls{TNE} profiles remain distinctly separated from \gls{TEQ} further into the boundary layer. This highlights the role of turbulence in delaying thermal relaxation compared to the laminar or transitional regimes. Overall, these results demonstrate that while velocity statistics conform to classical scaling once turbulence is established, thermal nonequilibrium effects persist much further downstream, underscoring the fundamentally different relaxation behaviour of velocity and temperature fields in turbulent boundary layers.


\begin{figure}[t]
    \centering
    \fontsize{9}{10}\selectfont
    \includegraphics[width=0.80\textwidth]{resources/figures/breakdown/result_dns_thermal_reynolds_stresses.pdf}
    \caption[Favre-Reynolds stress components in semilocal wall units for Mach 6, $T_w=1800$ K.]{Favre-Reynolds stress components and root-mean-square fluctuations plotted against semilocal wall coordinate $y^*$ for Mach 6, $T_w=1800$ K. Left column shows $\overline{\rho}\,\widetilde{u''u''}$, $\overline{\rho}\,\widetilde{v''v''}$, $\overline{\rho}\,\widetilde{w''w''}$, and $-\overline{\rho}\,\widetilde{u''v''}$ normalised by wall shear stress. Right column shows corresponding RMS velocity fluctuations. Comparisons between \gls{CPG}, \gls{TEQ}, and \gls{TNE} formulations at multiple streamwise stations.}
    \label{figure:thermal_reynolds_stresses}
\end{figure}
% \subsubsection{Favre averaged Turbulence stresses}

% The Turbulence stress components at $x/\delta^{*}_{0}=654$ and $740$, extracted from Favre-averaged statistics and normalised by the mean wall shear stress $\tau_{w}$, exhibit the canonical ordering $\overline{\rho}\,\widetilde{u''u''} > \overline{\rho}\,\widetilde{v''v''} \sim \overline{\rho}\,\widetilde{w''w''}$, as shown in \Cref{figure:thermal_reynolds_stresses}. The streamwise component reaches $\overline{\rho}\,\widetilde{u''u''}/\tau_{w}\approx 10.5$ at $y^{*}\approx 30$, whilst $\overline{\rho}\,\widetilde{v''v''}$ and $\overline{\rho}\,\widetilde{w''w''}$ attain peak values of approximately $3.5$ and $1.0$–$1.3$ respectively. These observations are broadly consistent with canonical \gls{DNS} of hypersonic turbulent boundary layers \citep{ZhangDuanChoudhari2018,XuEtal2021}. Critically, across all stress components at both stations, the \gls{CPG}, \gls{TEQ}, and \gls{TNE} closures collapse onto nearly indistinguishable profiles. This collapse demonstrates that once turbulence is fully developed, the inclusion of vibrational nonequilibrium effects produces negligible changes to the Reynolds stress structure, a result that reflects the differing transition Reynolds numbers across the three thermal models rather than any fundamental alteration to the turbulent dynamics. Furthermore, these comparisons reveal trends consistent with the findings of \cite{DiRenzo2021} and \cite{Williamsetal2025}, confirming that turbulence drives the nonequilibrium effects rather than the reverse, and thermal excitations are only important for the transition process. 

\subsubsection{Favre-averaged turbulence stresses}

The turbulence stress components at $x/\delta^*_0 = 654$ and $740$, extracted from Favre-averaged statistics and normalised by the mean wall shear stress $\tau_w$, exhibit the canonical ordering $\overline{\rho}\,\widetilde{u''u''} > \overline{\rho}\,\widetilde{v''v''} \sim \overline{\rho}\,\widetilde{w''w''}$ as shown in \Cref{figure:thermal_reynolds_stresses}. The streamwise component reaches $\overline{\rho}\,\widetilde{u''u''}/\tau_w \approx 10.5$ at $y^* \approx 30$, whilst $\overline{\rho}\,\widetilde{v''v''}$ and $\overline{\rho}\,\widetilde{w''w''}$ attain peak values of approximately $3.5$ and $1.0$--$1.3$ respectively, consistent with canonical DNS of hypersonic turbulent boundary layers~\citep{ZhangDuanChoudhari2018,XuEtal2021}.

Critically, across all stress components at both stations, the \gls{CPG}, \gls{TEQ}, and \gls{TNE} closures collapse onto nearly indistinguishable profiles. This collapse demonstrates that once turbulence is fully developed, vibrational nonequilibrium produces negligible changes to Reynolds stress structure. The result reflects differing transition Reynolds numbers across the three thermal models rather than any fundamental alteration to turbulent dynamics. These findings are consistent with recent studies~\citep{DiRenzo2021,Williamsetal2025}, confirming that turbulence drives nonequilibrium effects rather than the reverse, with thermal excitations important only for the transition process.


% \begin{figure}[t]
%     \centering
%     \fontsize{9}{10}\selectfont
%     \includegraphics[width=0.88\textwidth]{resources/figures/breakdown/result_dns_thermal_rhop_rhoT.pdf}
%     \caption[Density and temperature statistics in semilocal wall units for Mach~6, $T_w=1800$~K.]%
%     {Profiles of Favre-averaged pressue (top row) and temperature (bottom row) together with their corresponding normalised root-mean-square fluctuations, plotted against the semilocal wall coordinate $y^{*}$ for the Mach~6, $T_w=1800$~K case. Blue lines correspond to the \gls{CPG} formulation and red lines to the \gls{TNE} formulation, with results shown at several streamwise stations and the final location emphasised.}
%     \label{figure:thermal_rho_T}
% \end{figure}

% \Cref{figure:thermal_rho_T} shows the development of mean and rms pressure and temperature statistics across the boundary layer. The mean pressure profiles (top left) vary only weakly in the wall-normal direction, with all three formulations remaining close to their edge values. Small differences appear in the recovery region, where the \gls{TNE} solution relaxes more gradually than the \gls{CPG} and \gls{TEQ} cases, reflecting the slower adjustment associated with vibrational energy modes. This behaviour is consistent with the trends observed in the mean temperature (bottom left), where the \gls{TNE} state remains elevated further into the boundary layer. The fluctuation levels (right panels) display more distinct separation. Within the turbulent region, the \gls{TNE} formulation produces systematically larger rms pressure and temperature fluctuations than both \gls{CPG} and \gls{TEQ}. The additional degrees of freedom introduced by vibrational nonequilibrium strengthen the coupling between the thermal and momentum fields, which increases compressible-mode activity and enhances turbulent mixing. As a consequence, both pressure and temperature variances remain higher throughout the log layer and into the outer region.

% The temperature results exhibit particularly slow relaxation~\Cref{figure:thermal_rho_T} (bottom row). Linear \gls{DNS} studies suggested that the downstream \gls{TNE} temperature would approach the \gls{TEQ} state more closely than the \gls{CPG} state, yet the fully turbulent simulation demonstrates a markedly different behaviour. Turbulence acts to delay thermal equilibration, and the \gls{TNE} mean temperature remains clearly separated from \gls{TEQ} well into the boundary layer. These observations indicate that velocity statistics may adjust rapidly to classical turbulent behaviour, while thermal nonequilibrium persists much further downstream. Vibrational nonequilibrium is promoted by the presence of turbulence \citep{FievetEtAl2019}, meaning that once the flow has transitioned, it tends to remain in a nonequilibrium state over extended distances. This highlights that vibrational nonequilibrium is particularly relevant during the transition process, when the flow does not reach equilibrium, and understanding its dynamics is therefore important.

\begin{figure}[t]
    \centering
    \fontsize{9}{10}\selectfont
    \includegraphics[width=0.80\textwidth]{resources/figures/breakdown/result_dns_thermal_rhop_rhoT.pdf}
    \caption[Pressure and temperature statistics in semilocal wall units for Mach 6, $T_w=1800$ K.]{Profiles of Favre-averaged pressure (top row) and temperature (bottom row) with corresponding normalised root-mean-square fluctuations plotted against semilocal wall coordinate $y^*$ for Mach 6, $T_w=1800$ K. Left column shows mean profiles, with right showing the RMS fluctuations. Legend: \gls{CPG} (\protect\cpgline), \gls{TEQ} (\protect\teqline), \gls{TNE} (\protect\tneline).}
    \label{figure:thermal_rho_T}
\end{figure}

The mean pressure and temperature profiles, along with their RMS fluctuations, are presented across the boundary layer. The mean pressure varies only weakly in the wall-normal direction, with all three formulations remaining close to their edge values. Small differences appear in the recovery region, where \gls{TNE} relaxes more gradually than \gls{CPG} and \gls{TEQ}, reflecting slower adjustment associated with vibrational energy modes. The mean temperature shows similar behaviour, with \gls{TNE} remaining elevated further into the boundary layer. The fluctuation levels display more distinct separation, with \gls{TNE} producing larger RMS pressure and temperature fluctuations than both \gls{CPG} and \gls{TEQ}. The additional degrees of freedom introduced by vibrational nonequilibrium strengthen coupling between thermal and momentum fields, increasing compressible-mode activity and enhancing turbulent mixing. Consequently, both pressure and temperature variances remain higher throughout the log layer and outer region.

The temperature results exhibit particularly slow relaxation. Whilst linear DNS studies suggest downstream \gls{TNE} temperature should approach \gls{TEQ} more closely than \gls{CPG}, fully turbulent simulation shows markedly different behaviour. Turbulence acts to delay thermal equilibration, with \gls{TNE} mean temperature remaining clearly separated from \gls{TEQ} well into the boundary layer. These observations indicate that velocity statistics adjust rapidly to classical turbulent behaviour whilst thermal nonequilibrium persists much further downstream. Vibrational nonequilibrium is promoted by turbulence~\citep{FievetEtAl2019}, meaning that once the flow transitions, it tends to remain nonequilibrium over extended distances. This underscores that vibrational nonequilibrium effects are particularly significant during transition when the flow remains far from equilibrium.




% The temperature results exhibit particularly slow relaxation. Linear \gls{DNS} studies suggested that the downstream \gls{TNE} temperature would draw closer to the \gls{TEQ} state than to \gls{CPG}. The fully turbulent simulation shows a different outcome. Turbulence delays thermal equilibration, and the \gls{TNE} mean temperature remains distinctly separated from \gls{TEQ} far into the boundary layer. These results indicate that velocity statistics may settle rapidly into classical turbulent behaviour, yet thermal nonequilibrium persists much further downstream. Overall, the pressure and temperature fields highlight the different relaxation characteristics of high-speed turbulent boundary layers and reinforce the need to model vibrational nonequilibrium with care.



% \Cref{figure:thermal_rho_T} presents the evolution of rms pressure and temperature statistics across the boundary layer. In the mean profiles, for all three cases, 

% both \gls{CPG} and \gls{TNE} display broadly similar behaviour, with the density and temperature ratios decreasing from elevated near-wall values towards the freestream. Subtle differences emerge downstream, where the \gls{TNE} solutions relax more gradually, reflecting the delayed equilibration associated with vibrational energy modes.

% More pronounced deviations are observed in the fluctuation levels (right panels). Within the turbulent region, the \gls{TNE} formulation consistently exhibits larger root-mean-square variations in both density and temperature compared to the \gls{CPG} case. This amplification arises from the additional degrees of freedom introduced by vibrational nonequilibrium, which enhance the coupling between thermal and momentum fluctuations. The resulting increase in density and temperature variances strengthens turbulent mixing and alters the redistribution of energy within the boundary layer. In the mean temperature profiles, the \gls{CPG} and \gls{TNE} cases exhibit similar near-wall behaviour, but the \gls{TNE} solution decays more gradually towards the freestream, reflecting delayed vibrational relaxation. These findings indicate that, even when mean profiles appear to converge, the turbulence structure remains sensitive to the inclusion of vibrational nonequilibrium, with direct consequences for predicting wall heat transfer and informing model development.



% \section{Grid Study of DNS and ILES}
% In the following section, there will be a grid study of ILES cases to DNS cases, 

% \begin{figure}[t]
%     \centering
%     \fontsize{9}{10}\selectfont
%     \includegraphics[width=0.85\textwidth]{resources/figures/test/result_wallquantities_cf_st.pdf}
%     \caption[Streamwise evolution of $C_f$ and $St$ for Mach~6, $T_w=1800$ K.]%
%     {Streamwise evolution of spatially temporally averaged $C_f$ and $St$ for the Mach~6, $T_w=1800$ K case. Results are shown for \gls{CPG}, \gls{TEQ}, and \gls{TNE} formulations, alongside laminar reference solutions and the Van Driest~II turbulent correlation.}
%     \label{figure:gridstudy_cfst}
% \end{figure}

% % \subsection{Transitional regime}

% For that reason, the thesis concludes that ILES cases, although not good at detecting the turbulence intensity, are good at getting the transition onset location prediction and the nonlinear growth of the 2D modes. 

% \section{Role of Mach number and wall thermal conditions}

% \subsubsection{Results}
% \subsection{Breakdown to turbulence}
% Turbulent mean and flow fluctuations
% \subsubsection{Thermal nonequilibrium effects}
% \subsubsection{Heat flux and Reynolds stresses}

% \section{Role of Mach number and wall thermal conditions}

% % \subsection{Results}
% \subsection{Global flow properties}
% \subsection{Boundary-layer profiles}
% % \subsection{Shock--boundary layer interaction}

% \subsection{Breakdown to turbulence}
% \subsubsection{Onset location and growth rates}
% \subsubsection{Mode interactions and energy transfer}
% \subsubsection{Thermal nonequilibrium effects}
% \subsubsection{Heat flux and Reynolds stresses}
% \subsubsection{Second-order statistics}
% % \subsubsection{Spectral analysis of turbulence}

% \newpage

% \section{Role of Mach number and wall thermal temperature}
% The combined influence of Mach number and wall temperature conditions is now assessed within the framework of thermal nonequilibrium. These parameters exert a strong control on the stability and transition characteristics of hypersonic boundary layers: increasing Mach number enhances compressibility effects and amplifies acoustic-trapped disturbances, while variation in wall temperature modifies the mean thermal gradients that drive instability growth. To isolate these effects, simulations are performed for \gls{TNE} flows at Mach~6 and Mach~8, each with both a cold wall at $T_w=1200$~K and a hot wall at $T_w=1800$~K. This configuration enables a systematic examination of how freestream Mach number and wall thermal state interact to influence instability amplification, breakdown mechanisms, and the resulting surface heating environment. The numerics for the cases is mention in \Cref{section:breakdown:numericalsetup}. The wall-normal disturbance frequencies are imposed with non-dimensional frequency $f$. For Mach~6 cases, $f = 0.15$ and $0.16$; Mach~8 cases employ $f = 0.10$ and $0.11$.

% \section{Role of Mach number and wall thermal temperature}
% The combined influence of Mach number and wall temperature conditions is now assessed within the framework of thermal nonequilibrium. These parameters exert strong control over the stability and transition characteristics of hypersonic boundary layers: increasing Mach number enhances compressibility effects and amplifies acoustic-trapped disturbances, whilst wall temperature variation modifies the mean thermal gradients that drive instability growth. To isolate these effects, simulations are performed for \gls{TNE} flows at Mach~6 and Mach~8, each with a cold wall at $T_w=1200$~K and a hot wall at $T_w=1800$~K. This configuration enables systematic examination of how freestream Mach number and wall temperature interact to influence instability amplification, breakdown mechanisms, and surface heating. The numerical setup is described in \Cref{section:breakdown:numericalsetup}.

% Disturbance frequencies are imposed in non-dimensional form as $f = 0.15$ and $0.16$ for Mach~6 cases, and $f = 0.10$ and $0.11$ for Mach~8 cases. These values were selected to capture the resonance peak region, corresponding to physical domain locations around $x/\delta^*_0 = 200$ --  $400$ where modal growth is most pronounced (see \Cref{section:comparisonacrossmach}). The $N$-factor analysis in~\Cref{chapter:disturbance_growth} showed that Mach~6 cold wall case exhibits the highest $N$-factors, with the hot wall showing reduced growth, whilst both Mach~8 cases display further suppression of amplification, with the hot wall configuration providing the most stabilisation. This ranking demonstrates how increased Mach number and elevated wall temperature both act to stabilise modal growth, motivating the detailed examination of transition mechanisms in the sections that follow. The following section extends this framework isolating the \gls{TNE} effects across the cases. 


% \section{Role of Mach number and wall thermal temperature}

% The combined influence of Mach number and wall temperature on stability and transition in hypersonic boundary layers is assessed within vibrational nonequilibrium. Increasing Mach number amplifies compressibility effects and acoustic-trapped disturbances, whilst wall temperature variation modifies the mean thermal gradients that drive instability growth. To isolate these effects, simulations are performed for \gls{TNE} flows at Mach~6 and Mach~8 with cold wall $T_w = 1200$~K and hot wall $T_w = 1800$~K conditions. This enables systematic examination of how freestream Mach number and wall temperature interact to influence instability amplification, breakdown mechanisms, and surface heating, with the numerical setup described in \Cref{section:breakdown:numericalsetup}.

% Disturbance frequencies are imposed as $f = 0.15$ and $0.16$ for Mach~6, and $f = 0.10$ and $0.11$ for Mach~8, selected to capture the resonance peak region at $x/\delta^*_0 \approx 200$--$400$ where modal growth is most pronounced (see \Cref{section:comparisonacrossmach}). The $N$-factor analysis in \Cref{chapter:disturbance_growth} reveals that the Mach~6 cold wall case exhibits the highest amplification, with the hot wall showing suppression, whilst both Mach~8 cases display further stabilisation, most pronounced for the hot wall configuration. This ranking demonstrates that increased Mach number and elevated wall temperature both stabilise modal growth. The following section examines transition mechanisms in detail, isolating \gls{TNE} effects across these cases.


% \subsection{Global mean flow and surface properties}

% The global characteristics of the mean flow and surface interactions are examined through integrated aerothermodynamic quantities. The skin friction coefficient $C_f$ and heat transfer coefficient $C_q$~\citep{Volpiani2021} provide quantitative measures of wall shear stress and surface heating respectively, serving as key indicators of boundary layer development and transition onset across the range of Mach numbers and thermal conditions investigated. Tabulated values of these coefficients at representative streamwise locations enable direct comparison of how thermochemical nonequilibrium effects influence aerothermodynamic performance across the cold and hot wall configurations.

% \section{Grid study for DNS and ILES cases}
% \label{section:gridstudy}

% A grid refinement study was performed to ensure that the simulations adequately resolve both the mean boundary layer and the dominant instability modes. Three progressively finer meshes were tested, with particular attention to the near-wall spacing, which was kept below $\Delta y^+ < 1$, and sufficient streamwise and spanwise resolution to capture the relevant scales of motion. The fine-resolution \gls{DNS} grids comprised approximately $5200 \times 600 \times 275$ nodes in the streamwise, wall-normal, and spanwise directions respectively. Key boundary layer quantities, such as displacement thickness $\delta^*$, momentum thickness $\theta$, and enthalpy thickness $\delta_h$, were found to converge with refinement, indicating grid independence. Furthermore, the growth rates of linear \gls{DNS} instability modes were consistent across the grids, and inspection of instantaneous flow fields confirmed the absence of numerical artefacts.

% \begin{figure}[t]
%     \centering
%     \fontsize{9}{10}\selectfont
%     \includegraphics[width=0.88\textwidth]{resources/figures/breakdown/result_wallquantities_cf_st_grid.pdf}
%     \caption{Streamwise evolution of spatially temporally averaged skin friction coefficient $C_f$ and Stanton number $St$. Panels show: (a) CPG, (b) TEQ, and (c) TNE formulations, with laminar and turbulent reference solutions indicated.}
%     \label{figure:gridstudy_cfst}
% \end{figure}

% In addition to the fine-resolution \gls{DNS} cases, a coarser grid strategy was adopted to evaluate the feasibility of \gls{ILES} as a lower-cost alternative for transition studies. \gls{ILES} relies on the dissipative properties of high-order numerical schemes to implicitly model subgrid-scale turbulence, without explicit subgrid models \cite{Poggie2016}. The \gls{ILES} grids comprised approximately $2300 \times 500 \times 165$ nodes, representing a reduction of roughly 56\% in total grid points compared to the \gls{DNS} configuration. The grid spacings for the \gls{ILES} simulations were deliberately relaxed compared to the \gls{DNS} grids, resulting in streamwise and spanwise resolutions on the order of $\Delta x^+ \approx 20\mbox{--}27$ and $\Delta z^+ \approx 14\mbox{--}18$ respectively, whilst maintaining adequate wall-normal resolution with $\Delta y^+_w < 1$ in the viscous sublayer. This coarser resolution substantially reduces computational cost whilst still capturing the dominant instability mechanisms.

\section{Grid study for DNS and ILES cases}
\label{section:gridstudy}

The fine-resolution \gls{DNS} grids presented in the preceding sections comprised approximately $5200 \times 600 \times 275$ nodes in the streamwise, wall-normal, and spanwise directions respectively, with near-wall spacing designed to resolve viscous sublayer dynamics and instability modes. The domain dimensions are $800\delta^*_0 \times 100\delta^*_0 \times 20\delta^*_0$ in the streamwise, wall-normal, and spanwise directions, yielding grid spacings of $\Delta x^+ \approx 10.66$, $\Delta y^+_w \approx 0.94$, and $\Delta z^+ \approx 4.95$ respectively.

\begin{figure}[t]
    \centering
    \fontsize{9}{10}\selectfont
    \includegraphics[width=0.88\textwidth]{resources/figures/breakdown/result_wallquantities_cf_st_grid.pdf}
    \caption[Skin friction coefficient and Stanton number evolution for DNS simulations.]{Streamwise evolution of skin friction coefficient $C_f$ and Stanton number $St$ for the fine-resolution \gls{DNS} cases. Panels show (a) \gls{CPG}, (b) \gls{TEQ}, and (c) \gls{TNE} formulations with laminar and turbulent reference solutions.}
    \label{figure:gridstudy_cfst}
\end{figure}

In addition to the fine-resolution \gls{DNS} cases, a coarser grid strategy was adopted to evaluate \gls{ILES} as a lower-cost alternative for transition studies. \gls{ILES} relies on the dissipative properties of high-order numerical schemes to implicitly model subgrid-scale turbulence without explicit subgrid models~\citep{Poggie2016}. The \gls{ILES} grids comprised approximately $2300 \times 500 \times 165$ nodes over the same domain dimensions. This yields streamwise and spanwise resolutions of $\Delta x^+ \approx 23.33$ and $\Delta z^+ \approx 12.12$ respectively. This coarser resolution substantially reduces computational cost, with \gls{ILES} simulations approximately 6.5 times less expensive than \gls{DNS} due to reduced grid requirements, whilst capturing the dominant instability mechanisms.

\begin{figure}[t]
    \centering
    \fontsize{9}{10}\selectfont
    \includegraphics[width=0.88\textwidth]{resources/figures/breakdown/result_dns_thermal_boundarylayergrowth_grid3x3.pdf}
    \caption[Boundary-layer growth for DNS and ILES cases across thermal models.]{Streamwise evolution of boundary-layer growth showing displacement thickness $\delta^*$, momentum thickness $\theta^*$, and enthalpy thickness $\delta_h$ for (a) \gls{CPG}, (b) \gls{TEQ}, and (c) \gls{TNE} formulations. Solid lines represent \gls{DNS}, markers represent \gls{ILES}.}
    \label{figure:gridstudy_boundarylayergrowth}
\end{figure}

\begin{table}[htb]
\centering
\caption{Grid resolution parameters for DNS and ILES cases at Mach~6, wall temperature $T_w = 1800$~K.}
\label{table:grid_parameters}
\begin{tabular}{llcccccc}
\toprule
\textbf{Model} & \textbf{Case} & $\Delta x^+$ & $\Delta y^+_w$ & $\Delta y^+_{\delta_{99}}$ & $\Delta z^+$ & $L_z^+$ \\
\midrule
\multirow{2}{*}{CPG} & DNS & 10.66 & 0.94 & 9.15 & 4.95 & 1385 \\
& ILES & 24.18 & 0.88 & 8.52 & 16.24 & 1546 \\
\midrule
\multirow{2}{*}{TEQ} & DNS & 8.35 & 0.73 & 6.89 & 3.88 & 1085 \\
& ILES & 23.46 & 0.92 & 7.18 & 14.67 & 1239 \\
\midrule
\multirow{2}{*}{TNE} & DNS & 10.97 & 0.96 & 8.27 & 5.09 & 1426 \\
& ILES & 26.52 & 0.84 & 9.74 & 17.43 & 1678 \\
\bottomrule
\end{tabular}
\end{table}



To assess the suitability of \gls{ILES} for the present transition study, comparisons were made between corresponding \gls{DNS} and \gls{ILES} cases across the three thermal models. \Cref{figure:gridstudy_cfst} presents the streamwise evolution of skin friction coefficient $C_f$ and Stanton number $St$. The location of the Mack-mode resonance remains consistent between \gls{DNS} and \gls{ILES}, and the transition onset location is preserved. Notably, \gls{TNE} exhibits the sharpest transition in skin friction, reaching the transition threshold significantly earlier than \gls{CPG} and \gls{TEQ}. The primary differences between \gls{DNS} and \gls{ILES} manifest in the turbulent intensity and the sharpness of transition, where \gls{ILES} exhibits slightly smoother gradients in the transition region. This behaviour is consistent with observations that WENO and TENO schemes produce comparable results between \gls{DNS} and \gls{ILES}~\citep{HamzehloEtAl2021}.



\Cref{figure:gridstudy_boundarylayergrowth} shows the streamwise evolution of boundary-layer integral parameters. The displacement and momentum thicknesses track closely between \gls{DNS} and \gls{ILES} across all three thermal models, with \gls{TEQ} showing the tightest agreement and \gls{TNE} demonstrating remarkably consistent behaviour and trends throughout the domain. This close correspondence extends to the enthalpy thickness, indicating that \gls{ILES} accurately captures the overall boundary-layer growth and the transition location across thermal models.

\begin{figure}[t]
    \centering
    \fontsize{9}{10}\selectfont
    \includegraphics[width=0.88\textwidth]{resources/figures/breakdown/result_dns_grid_study_upH_Mt_k.pdf}
    \caption[Turbulence statistics comparison between DNS and ILES.]{Streamwise evolution of wall-normal velocity profiles $u^+$, turbulent Mach number $M_t$, and turbulent kinetic energy $k$ for (a) \gls{CPG}, (b) \gls{TEQ}, and (c) \gls{TNE} formulations. Solid lines represent \gls{DNS}, markers represent \gls{ILES}.}
    \label{figure:grid_study_upH_Mt_k}
\end{figure}

In the fully turbulent region, \Cref{figure:grid_study_upH_Mt_k} reveals minor differences in turbulence statistics. The \gls{ILES} velocity profiles $u^+$ follow the viscous sublayer and log law closely with \gls{DNS}, but begin to diverge in the buffer region and continue diverging throughout the log layer. The turbulent Mach number $M_t$ and turbulent kinetic energy $k$ (defined analogously to the disturbance kinetic energy in equation \Cref{equation:tke}
) show similarly modest overprediction. These differences arise from the coarser grid resolution underresolving fine-scale coherent structures in the turbulent boundary layer. However, the overall agreement is remarkably good considering the substantial reduction in computational cost, with \gls{ILES} approximately 6.4 times less expensive than \gls{DNS}. Consequently, whilst \gls{ILES} shows small errors in fully turbulent statistics, it provides a cost-effective pathway for systematic comparisons of transition location and early transition dynamics across multiple thermal models and flow conditions.



% To assess the suitability of \gls{ILES} for the present transition study, comparisons were made between corresponding \gls{DNS} and \gls{ILES} cases across the three thermal models. \Cref{figure:gridstudy_cfst} presents the streamwise evolution of skin friction coefficient $C_f$ and Stanton number $St$. The location of the Mack-mode resonance remains consistent between \gls{DNS} and \gls{ILES}, and the transition onset location is preserved within acceptable tolerance. The primary differences manifest in the turbulent intensity and the sharpness of transition, where \gls{ILES} exhibits slightly smoother gradients in the transition region. This behaviour is consistent with observations that WENO and TENO schemes produce comparable results between \gls{DNS} and \gls{ILES}~\citep{HamzehloEtAl2021}. Consequently, \gls{ILES} offers a computationally efficient alternative for transition studies and systematic comparisons across a wide range of conditions, Mach numbers, and wall temperatures.

% \begin{figure}[t]
%     \centering
%     \fontsize{9}{10}\selectfont
%     \includegraphics[width=0.88\textwidth]{resources/figures/breakdown/result_dns_thermal_boundarylayergrowth_grid3x3.pdf}
%     \caption{Streamwise evolution of spatially temporally averaged boundary layer properties. Columns show: (a) CPG, (b) TEQ, and (c) TNE formulations.}
%     \label{figure:gridstudy_boundarylayergrowth}
% \end{figure}


% However, \Cref{figure:gridstudy_boundarylayergrowth} reveals important limitations when assessing wall-bounded turbulence statistics. The displacement thickness $\delta^*$, momentum thickness $\theta$, and enthalpy thickness $\delta_h$ show close agreement between \gls{DNS} and \gls{ILES} in the transient nonlinear and early breakdown regions, where instability waves interact and secondary instabilities develop. In contrast, substantial deviations emerge in the fully turbulent region, where the coarser \gls{ILES} grid underresolves the turbulence structure and produces underpredicted skin friction. This limitation is well understood to arise from insufficient spatial resolution of fine-scale coherent structures in the turbulent boundary layer. Consequently, whilst \gls{ILES} is not suitable for detailed predictions of turbulent skin friction and wall-bounded turbulence statistics, it provides a viable and cost-effective pathway for systematic comparisons of transition location and early transition dynamics across multiple thermal models and flow conditions.

% \begin{figure}[t]
%     \centering
%     \fontsize{9}{10}\selectfont
%     \includegraphics[width=0.88\textwidth]{resources/figures/breakdown/result_dns_grid_study_upH_Mt_k.pdf}
%     \caption{Streamwise evolution of spatially temporally averaged boundary layer properties. Columns show: (a) CPG, (b) TEQ, and (c) TNE formulations.}
%     \label{figure:grid_study_upH_Mt_k}
% \end{figure}








% \section{Role of Mach number and wall thermal temperature}
% The combined influence of Mach number and wall temperature on stability and transition in hypersonic boundary layers is assessed within vibrational nonequilibrium. Increasing Mach number amplifies compressibility effects and are shown to dampen acoustic-trapped disturbances (\Cref{section:comparisonacrossmach},\citep{BitterShepherd2015}), whilst wall temperature variation modifies the mean thermal gradients that drive instability growth. Simulations are performed for \gls{TNE} flows at Mach~6 and Mach~8 with cold wall $T_w = 1200$~K and hot wall $T_w = 1800$~K conditions (numerical setup in \Cref{section:breakdown:numericalsetup}).

% Disturbance frequencies are imposed as $f = 0.15$ and $0.16$ for Mach~6, and $f = 0.10$ and $0.11$ for Mach~8, selected to capture the resonance peak region at $x/\delta^*_0 \approx 200$--$400$ where modal growth is most pronounced (see \Cref{section:comparisonacrossmach}), and trigger transition within the domain. The $N$-factor analysis in \Cref{chapter:disturbance_growth} reveals that the Mach~6 cold wall case exhibits the highest amplification, with the hot wall showing suppression, whilst both Mach~8 cases display further stabilisation. This demonstrates that increased Mach number and elevated wall temperature both stabilise modal growth.

\section{Role of Mach number and wall temperature}

The combined influence of Mach number and wall temperature on stability and transition in hypersonic boundary layers is assessed within vibrational nonequilibrium. Increasing Mach number amplifies compressibility effects that dampen acoustic-trapped disturbances~\citep{BitterShepherd2015}, whilst wall temperature variation modifies the mean thermal gradients driving instability growth. Simulations are performed for \gls{TNE} flows at Mach 6 and Mach 8 with cold wall $T_w = 1200~\text{K}$ and hot wall $T_w = 1800~\text{K}$ conditions using \gls{ILES}. The computational domains span $800\delta^*_0$ for Mach 6 and $1600\delta^*_0$ for Mach 8 to accommodate different instability wavelengths and growth rates. Grid resolutions are specified in \Cref{table:domain_grid_iles}, with streamwise spacing adjusted to maintain consistent wall-normal and spanwise resolution across Mach numbers.

Disturbance frequencies are imposed as $f = 0.15$ and $0.16$ for Mach 6, and $f = 0.10$ and $0.11$ for Mach 8, selected to capture the resonance peak region at $x/\delta^*_0 \approx 200$--$400$ where modal growth is most pronounced and trigger transition within the domain. The $N$-factor analysis in \Cref{chapter:disturbance_growth} reveals that the Mach 6 cold wall case exhibits the highest amplification, with the hot wall showing suppression, whilst both Mach 8 cases display further stabilisation. This demonstrates that increased Mach number and elevated wall temperature both stabilise modal growth. Vibrational nonequilibrium effects persist across all Mach numbers and wall temperatures examined.

% \begin{figure}[t]
%     \centering
%     \fontsize{9}{10}\selectfont
%     \includegraphics[width=0.95\textwidth]{resources/figures/breakdown/result_wallquantities_cf_st_tne.pdf}
%     \caption[Skin friction coefficient evolution for TNE cases across Mach 6 and 8.]{Streamwise evolution of skin friction coefficient $C_f$ and Stanton number $St$ for hypersonic boundary layers with varying Mach number and wall temperature. Panels show (a) Mach 6, $T_w=1200$ K; (b) Mach 6, $T_w=1800$ K; (c) Mach 8, $T_w=1200$ K; and (d) Mach 8, $T_w=1800$ K for \gls{TNE} formulations with laminar and turbulent reference solutions.}
%     \label{figure:mach_cf}
% \end{figure}


\subsection{Global mean flow and surface properties}

The skin friction coefficient $C_f$, Stanton number $St$, and normalised heat transfer coefficient $C_{q,w} = q_w / (\rho_e u_e^3)$ provide quantitative measures of wall shear stress and surface heating across all Mach numbers and thermal conditions~\citep{Volpiani2021}. \Cref{figure:mach_cf_st} presents the skin friction and Stanton number evolution across all cases. A pronounced overshoot in $C_f$ occurs upon transition onset across all cases. The resonance peak is pronounced in the Mach 6 cases, consistent with earlier DNS studies, but considerably subdued in the Mach 8 cases. This attenuation reflects the lower $N$-factor response at higher Mach number, with maximum $N$-factors of 8 and 6 for Mach 6 cases compared to 4 and 2 for Mach 8 cases. The transition sequence follows the expected progression, with the Mach 6 cold wall case transitioning earliest and the Mach 8 hot wall case exhibiting the most extended transition process. Notably, despite employing the same grid resolution as the Mach 6 cases, the Mach 8 simulations exhibit a DNS-level spike in skin friction during transition, indicating that higher compressibility effects at elevated Mach number produce sharper transition characteristics even on coarser grids. The Stanton number follows similar trends to skin friction, with pronounced peaks in the Mach 6 cases and more gradual evolution in the Mach 8 cases.

\begin{figure}[t]
    \centering
    \fontsize{9}{10}\selectfont
    \includegraphics[width=0.95\textwidth]{resources/figures/breakdown/result_wallquantities_cf_st_tne.pdf}
    \caption[Skin friction and Stanton number evolution for TNE cases across Mach and wall temperature.]{Streamwise evolution of skin friction coefficient $C_f$ and Stanton number $St$ for \gls{TNE} formulations showing (a) Mach 6, $T_w=1200$ K; (b) Mach 6, $T_w=1800$ K; (c) Mach 8, $T_w=1200$ K; and (d) Mach 8, $T_w=1800$ K with laminar and turbulent reference solutions.}
    \label{figure:mach_cf_st}
\end{figure}

\begin{figure}[h]
    \centering
    \fontsize{9}{10}\selectfont
    \includegraphics[width=0.75\textwidth]{resources/figures/breakdown/result_wallquantities_cq_tne.pdf}
    \caption[Heat flux components for TNE cases across Mach and wall temperature.]{Streamwise evolution of translational--rotational ($C_{q,\mathrm{tr}}$, black) and vibrational ($C_{q,\mathrm{v}}$, red, scaled $\times 10$) heat flux coefficients for \gls{TNE} formulations showing (a) Mach 6, $T_w=1200$ K; (b) Mach 6, $T_w=1800$ K; (c) Mach 8, $T_w=1200$ K; and (d) Mach 8, $T_w=1800$ K.}
    \label{figure:cqplot}
\end{figure}

To assess the influence of vibrational nonequilibrium, the translational--rotational and vibrational heat flux components are examined. The net total heat flux is represented by the sum of $C_{q,\mathrm{tr}}$ and $C_{q,\mathrm{v}}$, normalised as $C_{q,w} = q_w / (\rho_e u_e^3)$ as shown in \Cref{figure:cqplot}, with $C_{q,\mathrm{v}}$ scaled by a factor of ten.

The translational--rotational and vibrational heat flux components exhibit strikingly different behaviours across the four cases. In the Mach 6 cold wall case, $C_{q,\mathrm{v}}$ remains suppressed during the laminar and linear growth regions, rising sharply coincident with secondary instability and transition. The Mach 6 hot wall case shows elevated vibrational heat flux throughout, with significant levels in the laminar region that amplify further during transition. At Mach 8, the cold wall case exhibits minimal vibrational contribution across the entire domain, whilst the hot wall case shows weak but persistent vibrational activity that gradually increases during the transition process. This dependence on wall temperature reveals that vibrational excitation is promoted by elevated temperatures, with cold walls suppressing vibrational modes until nonlinear breakdown occurs.

% \subsection{Global mean flow and surface properties}

% The skin friction coefficient $C_f$, $St$ and normalised heat transfer coefficient $C_{q,w} = q_w / (\rho_e u_e^3)$ provide quantitative measures of wall shear stress and surface heating across all Mach numbers and thermal conditions~\citep{Volpiani2021}. \Cref{figure:mach_cf} presents the skin friction and Stanton number evolution across all cases. A pronounced overshoot in $C_f$ occurs upon transition onset across all cases. The resonance peak is pronounced in the Mach 6 cases, consistent with earlier DNS studies, but considerably subdued in the Mach 8 cases. This attenuation reflects the lower $N$-factor response at higher Mach number, with maximum $N$-factors of 8 and 6 for Mach 6 cases compared to 4 and 2 for Mach 8 cases. The transition sequence follows the expected progression, with the Mach 6 cold wall case transitioning earliest and the Mach 8 hot wall case exhibiting the most extended transition process.

% \begin{figure}[h]
%     \centering
%     \fontsize{9}{10}\selectfont
%     \includegraphics[width=0.75\textwidth]{resources/figures/breakdown/result_wallquantities_cq_tne.pdf}
%     \caption[Heat flux components for TNE cases across Mach and wall temperature conditions.]{Streamwise evolution of translational--rotational ($C_{q,\mathrm{tr}}$, black) and vibrational ($C_{q,\mathrm{v}}$, red, scaled $\times 10$) heat flux coefficients for \gls{TNE} formulations showing (a) M6Tw12, (b) M6Tw18, (c) M8Tw12, and (d) M8Tw18.}
%     \label{figure:cqplot}
% \end{figure}

% To assess the influence of vibrational nonequilibrium, the wall heat fluxes for translational--rotational and vibrational modes are examined, normalised as $C_{q,w} = q_w / (\rho_e u_e^3)$. The net total heat flux is represented by the sum of the translational--rotational component $C_{q,\mathrm{tr}}$ and the vibrational component $C_{q,\mathrm{v}}$. \Cref{figure:cqplot} presents the normalised heat flux evolution across all cases, with $C_{q,\mathrm{v}}$ scaled by a factor of ten to facilitate visual comparison.

% The translational--rotational and vibrational heat flux components exhibit markedly different behaviours across the flow development. Whilst both modes follow trends qualitatively similar to the skin friction coefficient, the vibrational contribution is substantially lower in magnitude and demonstrates a pronounced dependence on wall temperature. For cold wall flows, $C_{q,v}$ remains approximately constant throughout the laminar region until the secondary instability, whereupon it increases rapidly in conjunction with transition onset, consistent with prior observations \citep{Passiatore2021}. Conversely, hot wall cases exhibit a distinctly different response, the vibrational mode responds not solely to the secondary instability but tracks the translational--rotational evolution throughout both laminar and turbulent regions, albeit with reduced amplitude. This contrasting behaviour reflects the competing effects of thermal conductivity ($\kappa$) and vibrational thermal conductivity ($\kappa_v$) across the range of wall temperatures examined.


% \begin{figure}[t]
%     \centering
%     \fontsize{9}{10}\selectfont
%     \includegraphics[width=0.95\textwidth]{resources/figures/breakdown/result_wallquantities_cf_st_tne.pdf}
%     \caption[Streamwise evolution of spatially temporally averaged skin friction coefficient $C_f$. Panels show: (a)~Mach~6, $T_w=1200$~K; (b)~Mach~6, $T_w=1800$~K; (c)~Mach~8, $T_w=1200$~K; and (d)~Mach~8, $T_w=1800$~K. Results are presented for TNE formulations, with laminar and turbulent reference solutions indicated.]%
%     {Streamwise evolution of spatially temporally averaged skin friction coefficient $C_f$ for hypersonic boundary layers with varying thermochemical treatments. Panels show: (a)~Mach~6, $T_w=1200$~K; (b)~Mach~6, $T_w=1800$~K; (c)~Mach~8, $T_w=1200$~K; and (d)~Mach~8, $T_w=1800$~K. Results are presented for\gls{TNE} formulations, with laminar and turbulent reference solutions indicated.}
%     \label{figure:mach_cf}
% \end{figure}
% \subsection{Global mean flow and surface properties}

% The skin friction coefficient $C_f$ and normalised heat transfer coefficient $C_{q,w} = q_w / (\rho_e u_{e}^{3})$ provide quantitative measures of wall shear stress and surface heating, serving as key indicators across all investigated Mach numbers and thermal conditions \citep{Volpiani2021}. \Cref{figure:mach_cf} presents the skin friction evolution across all cases. A notable feature common to all cases is the pronounced overshoot in $C_f$ observed upon the onset of turbulence. This transient overshoot. Whilst a pronounced peak is evident in the Mach~6 cases, consistent with earlier \gls{DNS} studies, the resonance peak is considerably subdued in the Mach~8 cases. This attenuation may be attributed to the relatively low $N$-factor response characteristic of these conditions. For reference, the maximum $N$-factors are 8 and 6 for the Mach~6 cases, compared to 4 and 2 for the Mach~8 cases. The transition sequence follows the expected progression, with the Mach~6 cold wall case transitioning to turbulence earliest, whilst the Mach~8 hot wall case exhibits the most extended transition process.



% \begin{figure}[h]
%     \centering
%     \fontsize{9}{10}\selectfont
%     \includegraphics[width=0.75\textwidth]{resources/figures/breakdown/result_wallquantities_cq_tne.pdf}
%     \caption[Streamwise evolution of spatially temporally averaged translational--rotational ($C_{q,tr}$, black) and vibrational ($C_{q,v}$, red, scaled $\times 10$) heat flux coefficients for the Mach~6, $T_w = 1800$ K case. Results are shown for TNE formulations, (a) M6Tw12, (b) M6Tw18, (c) M8Tw12, (d) M8Tw18]%
%     {Streamwise evolution of spatially temporally averaged translational--rotational ($C_{q,tr}$, black) and vibrational ($C_{q,v}$, red, scaled $\times 10$) heat flux coefficients for the Mach~6, $T_w = 1800$ K case. Results are shown for \gls{TNE} formulations, (a) M6Tw12, (b) M6Tw18, (c) M8Tw12, (d) M8Tw18.}
%     \label{figure:cqplot}
% \end{figure}

% To properly assess the influence of vibrational nonequilibrium, the wall heat fluxes for the translational--rotational and vibrational modes are examined. These are normalised using the definition $C_{q,w} = q_w / (\rho_e u_{e}^{3})$, with the net total heat flux represented by the sum of the translational--rotational component $C_{q,tr}$ and the vibrational component $C_{q,v}$. \Cref{figure:cqplot} presents the normalised heat flux evolution across all cases, with $C_{q,v}$ scaled by a factor of ten to facilitate visual comparison within the domain. 

% Marked differences emerge between the two modes across the flow development. Both components exhibit trends qualitatively similar to the skin friction coefficient in all cases; however, the vibrational contribution is substantially lower in magnitude and demonstrates distinctly different behaviour between cold wall and hot wall configurations. For cold wall flows, $C_{q,v}$ remains relatively constant throughout the laminar region up to the secondary instability, beyond which it undergoes a rapid increase consistent with the onset of transition, as reported by \cite{Passiatore2021}. In contrast, hot wall cases exhibit a qualitatively different response, wherein the vibrational energy mode responds not only to the secondary mode instability but also follows the translational--rotational heat flux evolution in both laminar and turbulent regions, albeit with reduced amplitude. The contrast differences show the effect of $\kappa$ and $\kappa_v$ at different wall temperatures. 








% \begin{figure}[h]
%     \centering
%     \fontsize{9}{10}\selectfont
%     \includegraphics[width=0.80\textwidth]{resources/figures/breakdown/result_wallquantities_cq_tne.pdf}
%     \caption[Streamwise evolution of $C_f$ and $St$ for Mach~6, $T_w=1800$ K.]%
%     {Streamwise evolution of spatially temporally averaged $C_f$ and $St$ for the Mach~6, $T_w=1800$ K case. Results are shown for \gls{CPG}, \gls{TEQ}, and \gls{TNE} formulations, alongside laminar reference solutions and the Van Driest~II turbulent correlation.}
%     \label{figure:cqplot}
% \end{figure}






% Understanding the breakdown process, which is significantly delayed in the Mach 8 cases, what causes this breakdown? 

% \subsection{Onset location and growth rates}
% \subsubsection{Boundary-layer profiles}

% \subsection{Mode interactions and energy transfer}

% Similar work of POD from last chapter, but with some extra results to show dominant modes of Mack 8 and hopefully answer the question of why there is a dip in Mach 6 cases and not the Mach 8 case. 

% \subsection{Breakdown to turbulence}





% \subsection{Streak breakdown}
% The main section to talk about I guess, Streak amplitude plot, breakdown to turbulence due to the prescence of streaks, \citep{Andersson2001} and \citep{Franko2013} have spoke  about this. \cite{BrandtBottaro2001} incompressible 26\% of freestream. 



% \subsection{Streak breakdown and transition to turbulence}

% Streak breakdown via secondary instability is a critical transition mechanism in boundary layers. Streaks are quasi-streamwise vortical structures generated through the lift-up mechanism \citep{Andersson2001}, which amplify until secondary instability triggers breakdown to turbulence. The critical amplitude for breakdown has been established at approximately $26\,\%$ of the freestream velocity in incompressible flows \citep{BrandtBottaro2001}, with this threshold remaining robust across compressible regimes \citep{Franko2013}.

% \Cref{figure:streakbreakdown} shows streak amplitude evolution across the four thermal nonequilibrium cases. At Mach~6, $T_w = 1200$ K, rapid amplification occurs and the secondary instability threshold is crossed early (approximately $32\,\%$ of freestream). Increasing the wall temperature to $1800$ K substantially delays streak growth, reflecting the stabilising effect of elevated wall temperature on secondary instability. At Mach~8, streak amplitudes remain heavily suppressed at both wall temperatures, indicating that compressibility and wall heating combine to inhibit secondary instability at this higher Mach number.


\begin{figure}[t]
    \centering
    \includegraphics[width=9.3cm]{resources/figures/breakdown/result_tne_amplitude_nfactor.pdf}
    \caption[Linear amplitude growth and $N$-factor evolution for ILES forcing frequencies across Mach and wall temperature conditions.]{Streamwise evolution of rms amplitude (left column) and cumulative $N$-factor (right column) for the prescribed forcing frequencies showing (a) Mach 6, $T_w=1200$ K ; (b) Mach 6, $T_w=1800$ K; (c) Mach 8, $T_w=1200$ K; and (d) Mach 8, $T_w=1800$ K  for \gls{TNE} conditions.}
    \label{figure:tne_amplitude_nfactor}
\end{figure}



\subsection{Modal growth}

One of the distinct features of the transition studies presented in this chapter have been the resonance peaks in the skin friction and Stanton number plot. The resonance peak was attributed to the initial growth of the 2D Mack mode~\citep{FrankoLele2013}, and the peak was observed in subsequent studies of hypersonic boundary layers~\citep{DiRenzo2021}, though absent in others. Both observations stem from adiabatic wall conditions. The prescribed forcing frequencies are selected to excite the mode at its point of maximum linear amplification, as identified in the $N$-factor analysis presented in \Cref{figure:tne_amplitude_nfactor}, with the peak location corresponding to where the $N$-factor reaches its maximum. The resonance peak prominence is strongly modulated by thermal conditions, appearing pronounced in the Mach 6 cases where $N$-factors reach 8 and 6 respectively, but considerably subdued in the Mach 8 cases where $N$-factors remain at 4 and 2. Notably, despite employing identical forcing frequency amplitudes across all cases, the resonance peaks are absent in the Mach 8 simulations, reflecting the substantially lower linear growth rates at elevated Mach number that fail to produce sufficient wall-quantity amplification to manifest as distinct peaks. Cold wall conditions enhance the peak through increased thermal gradients driving instability growth, whilst elevated wall temperature and higher Mach number both suppress linear amplification. The resonance peaks transition smoothly into nonlinear breakdown as secondary instabilities emerge, reflecting the natural progression from linear to nonlinear dynamics.


\begin{figure}[H]
    \centering
    \includegraphics[width=22.3cm,angle=-90]{resources/figures/breakdown/result_iles_thermal_contour_pod_mode1.pdf}
    \caption[Leading POD mode of pressure fluctuations across Mach and wall temperature conditions.]{Streamwise evolution of the first \gls{POD} of pressure fluctuations in the spanwise direction for \gls{TNE} formulations showing (a) Mach 6, $T_w=1200$ K; (b) Mach 6, $T_w=1800$ K; (c) Mach 8, $T_w=1200$ K; and (d) Mach 8, $T_w=1800$ K.}
    \label{figure:iles_pod_mode1}
\end{figure}

% To further isolate the dominant coherent structures driving transition across the parametric range, proper orthogonal decomposition was applied to pressure fluctuations in the spanwise direction. \Cref{figure:iles_pod_mode1} presents the first POD mode of pressure fluctuations across all four thermal conditions. The Mach 6 cases display diffuse, irregular patterns that rapidly break down into small-scale fluctuations, reflecting the vigorous nonlinear interactions characteristic of nonlinear growth. In contrast, the Mach 8 cases exhibit highly organised, periodic spanwise structures that persist throughout the modal amplification regime, indicating a strong modulation by the forcing frequency. The spanwise wavelengths evident in the \gls{POD} mode align with the imposed forcing wavelengths, confirming that the dominant pressure structures are directly driven by the prescribed disturbances.

To further isolate the dominant coherent structures driving transition across the parametric range, proper orthogonal decomposition was applied to pressure fluctuations in the spanwise direction. \Cref{figure:iles_pod_mode1} presents the first POD mode of pressure fluctuations across all four thermal conditions. The Mach 6 cases display diffuse, irregular patterns that rapidly break down into small-scale fluctuations, reflecting the vigorous nonlinear interactions characteristic of nonlinear growth. In contrast, the Mach 8 cases exhibit highly organised, periodic spanwise structures that persist throughout the modal amplification regime, indicating a strong modulation by the forcing frequency. The spanwise wavelengths evident in the \gls{POD} mode align with the imposed forcing wavelengths, confirming that the dominant pressure structures are directly driven by the prescribed disturbances.

The presence of the resonance peak in the Mach 6 cases and its absence in the Mach 8 cases reflects the competition between linear amplification and nonlinear saturation. At Mach 6, the higher linear growth rates allow the forced mode to reach substantial amplitude over a distinct spatial region before secondary instabilities dominate, producing the pronounced resonance peak. At Mach 8, the suppressed linear growth rates prevent the primary mode from achieving sufficient amplitude to generate a detectable peak, with transition driven instead by a more gradual energy cascade directly into nonlinear breakdown.

% \begin{figure}[H]
%     \centering
%     \includegraphics[width=22.3cm,angle=-90]{resources/figures/breakdown/result_iles_thermal_contour_pod_mode1.pdf}
%     \caption[Time-averaged wall Stanton number ($St$) during boundary layer transition for (a) CPG, (b) TEQ, and (c) TNE models.]%
%     {Time-averaged wall Stanton number ($St$) during boundary layer transition for (a) \gls{CPG}, (b) \gls{TEQ}, and (c) \gls{TNE} models.}
%     \label{figure:thermal_spanwise_st}
% \end{figure}

\subsection{Streak breakdown}

Streak breakdown via secondary instability is a critical transition mechanism in boundary layers. Streaks are quasi-streamwise vortical structures generated through the lift-up mechanism, which amplify until secondary instability triggers breakdown to turbulence. In incompressible flows, secondary instabilities were characterised using Floquet theory~\citep{Andersson2001}, revealing sinuous and varicose modes with critical amplitudes of $26\,\%$ and $37\,\%$ of freestream velocity respectively. Brandt and Bottaro~\citep{BrandtBottaro2001} extended this work to include nonlinear saturation of streaks and full breakdown to turbulence through forced secondary instability. However, the critical amplitude threshold for compressible flows remains poorly established, with no dominant mechanism or commonly accepted streak amplitude percentage identified in high-speed simulations~\citep{Franko2013}.

Streak amplitude is computed from the velocity field as the maximum spanwise variation of streamwise velocity, normalised by the freestream velocity. For each streamwise location and wall-normal position, the spanwise deviation is calculated as
\begin{equation}
A_s(x) = \frac{\max_z u(x, y, z) - \min_z u(x, y, z)}{u_e},
\end{equation}
where the streak amplitude at each streamwise station is then defined as the maximum value across the wall-normal direction.



\begin{figure}[H]
    \centering
    \includegraphics[width=20.3cm,angle=-90]{resources/figures/breakdown/result_iles_streak_contour.pdf}
    \caption[Streak amplitude evolution across Mach and wall temperature conditions.]{Streamwise evolution of streak amplitude $A_s/u_e$ in the spanwise direction for \gls{TNE} formulations showing (a) Mach 6, $T_w=1200$ K; (b) Mach 6, $T_w=1800$ K; (c) Mach 8, $T_w=1200$ K; and (d) Mach 8, $T_w=1800$ K. }
    \label{figure:iles_streak_amplitude}
\end{figure}

The development of sinusoidal streak structures as mentioned in \citep{Andersson2001} is evident across all thermal conditions shown in \Cref{figure:iles_streak_amplitude}. The Mach 6 cases exhibit rapid streak formation and amplification beginning in the secondary instability region, with coherent spanwise-periodic structures emerging around $x/\delta^*_0 \approx 300$ and growing to substantial amplitudes approaching the breakdown region. The Mach 8 cases show more gradual streak development with lower peak amplitudes, consistent with the weaker secondary instability growth at higher Mach number. Streak amplitude increases systematically through the transition zone, with the most vigorous growth occurring in the Mach 6 cold wall case where linear instability is most pronounced. The sinusoidal organisation reflects the forcing-induced modulation of the flow, demonstrating that secondary instabilities preferentially amplify structures aligned with the imposed spanwise wavelengths.

\begin{figure}[t]
    \centering
    \fontsize{9}{10}\selectfont
    \includegraphics[width=0.90\textwidth]{resources/figures/breakdown/result_iles_streak_amplitude.pdf}
    \caption[Streamwise evolution of spatially temporally averaged translational--rotational ($C_{q,tr}$, black) and vibrational ($C_{q,v}$, red, scaled $\times 10$) heat flux coefficients for the Mach~6, $T_w = 1800$ K case. Results are shown for TNE formulations, (a) M6Tw12, (b) M6Tw18, (c) M8Tw12, (d) M8Tw18]%
    {Streak amplitude growth in the streamwise direciton. Results are shown for \gls{TNE} studies.}
    \label{figure:streakbreakdown}
\end{figure}


This metric quantifies the strength of quasi-streamwise vortical structures and provides a direct measure of secondary instability susceptibility. The streak amplitude evolution across the four thermal nonequilibrium cases is shown in \Cref{figure:streakbreakdown}. At Mach 6, $T_w = 1200$ K, rapid amplification occurs with the secondary instability threshold crossed early at approximately $32\,\%$ of freestream velocity. Increasing wall temperature to $1800$ K delays streak growth, reflecting the stabilising effect of elevated wall temperature on secondary instability. At Mach 8, streak amplitudes remain suppressed at both wall temperatures, indicating that compressibility and wall heating combine to inhibit secondary instability at higher Mach number. The thermal nonequilibrium effect manifests through altered thermodynamic properties governing secondary instability growth rates. Energy partitioning between translational and vibrational modes introduces a stabilising influence on streak amplification, consistent with the delayed transition observed in \gls{TNE} mean flow statistics. These results emphasise the importance of accounting for compressibility, wall effects, and thermal nonequilibrium when predicting transition in hypersonic flows.


% The thermal nonequilibrium effect manifests through altered thermodynamic properties governing secondary instability growth rates. The energy partitioning between translational and vibrational modes introduces a stabilising influence on streak amplification, consistent with the delayed transition observed in TNE mean flow statistics. These results underscore the necessity of accounting for compressibility, wall effects, and thermal nonequilibrium when predicting transition in hypersonic flows.




% \newpage

% \subsection{Thermal nonequilibrium effects}
% This is where we introduce the thermal nonequilibrium effects from statistics to how the delay in transition happens, 

\subsection{Thermal nonequilibrium effects}
% In this section, vibrational nonequilibrium effects across the four thermal conditions are analysed to elucidate how Mach number and wall temperature modulate transition behaviour in thermochemically reactive flows. Figure~\ref{figure:contour_TmTv} presents contours of vibrational temperature deviation $T_v - T$ normalised by freestream temperature throughout the domain, with non-unity aspect ratio to reveal regions of thermal nonequilibrium. The vibrational nonequilibrium zones are most prominent in the transition region where secondary instabilities drive rapid energy redistribution between translational and vibrational modes. In the fully turbulent region downstream, the nonequilibrium zones persist with similar magnitude, indicating that vibrational relaxation effects remain dynamically significant even after the transition to turbulence is complete. The spatial extent and intensity of nonequilibrium vary systematically with Mach number and wall temperature, with elevated Mach number and cold walls suppressing vibrational excitation whilst lower Mach number and hot walls promote stronger thermal nonequilibrium effects.

% In this section, vibrational nonequilibrium effects across the four thermal conditions are analysed to elucidate how Mach number and wall temperature modulate transition behaviour in hypersonic flows. Figure~\ref{figure:contour_TmTv} presents contours of time-averaged vibrational temperature deviation $T_v - T$ normalised by freestream temperature throughout the domain, with non-unity aspect ratio to reveal regions of thermal nonequilibrium. The maximum vibrational nonequilibrium occurs within the transition region across all cases, where secondary instabilities drive rapid energy redistribution between translational and vibrational modes. In the fully turbulent region downstream, the nonequilibrium zones persist with reduced magnitude, indicating that vibrational relaxation effects become less dynamically significant once transition is complete.  The spatial extent and intensity of nonequilibrium vary systematically with Mach number and wall temperature, with elevated Mach number and cold walls suppressing vibrational excitation whilst lower Mach number and hot walls promoting stronger thermal nonequilibrium effects. 

In this section, vibrational nonequilibrium effects across the four thermal conditions are analysed to understand how Mach number and wall temperature modulate transition behaviour in hypersonic flows. Figure~\ref{figure:contour_TmTv} presents contours of time- and spanwise-averaged vibrational temperature deviation $\overline{T_v - T}$ normalised by freestream temperature throughout the domain, with non-unity aspect ratio to reveal regions of thermal nonequilibrium. The maximum vibrational nonequilibrium occurs within the transition region across all cases, where secondary instabilities drive rapid energy redistribution between translational and vibrational modes. In the fully turbulent region downstream, the nonequilibrium zones persist with reduced magnitude, indicating that vibrational relaxation effects become less dynamically significant once transition is complete. The spatial extent and intensity of nonequilibrium vary systematically with Mach number and wall temperature, with elevated Mach number and cold walls suppressing vibrational excitation whilst lower Mach number and hot walls promoting stronger thermal nonequilibrium effects.



These observations are consistent with the findings of~\cite{FievetEtAl2019} and~\cite{Passiatore2021}, who identified that over-excited vibrational states develop in regions where turbulent intensity decreases whilst vibrational energy relaxes slowly. However, the present results show that the transition process itself drives the flow into thermal nonequilibrium through the intense energy redistribution accompanying secondary instability growth. The concentration of nonequilibrium in the transition region demonstrates that the instability-driven acceleration and deceleration of fluid parcels creates thermodynamic conditions where the vibrational relaxation timescale becomes comparable to the hydrodynamic timescale, preventing rapid equilibration and generating over-excited vibrational states. To quantify this balance and understand how the competing timescales determine the strength of TNE effects across the parametric range examined, a Damköhler number analysis based on the viscous timescale is performed to characterise the interaction between vibrational relaxation and instability growth.

\begin{figure}[t]
    \centering
    \fontsize{9}{10}\selectfont
    \includegraphics[width=0.80\textwidth]{resources/figures/breakdown/result_iles_contour_TmTv_fulldomain.pdf}
    \caption[Streamwise evolution of time averaged temperature difference normalised by freestream conditions for TNE formulations showing (a) Mach 6, $T_w=1200$ K; (b) Mach 6, $T_w=1800$ K; (c) Mach 8, $T_w=1200$ K; and (d) Mach 8, $T_w=1800$ K.]{Streamwise evolution of mean temperature field normalised by freestream conditions for \gls{TNE} formulations showing (a) Mach 6, $T_w=1200$ K; (b) Mach 6, $T_w=1800$ K; (c) Mach 8, $T_w=1200$ K; and (d) Mach 8, $T_w=1800$ K.}
    \label{figure:contour_TmTv}
\end{figure}

The streamwise evolution of the Damk\"{o}hler number at five characteristic stations is presented in \ref{figure:damkohlernumber}. The Damk\"{o}hler number, defined as the ratio of the viscous timescale $\delta^{*}/u_{\tau}$ to the vibrational relaxation timescale, quantifies the relative importance of instability-driven energy redistribution compared to molecular relaxation. Across all thermal conditions, a consistent trend emerges: the Damk\"{o}hler number peaks are progressively confined towards the wall during the secondary instability region. This wall-confinement indicates that vibrational excitation localises where the viscous timescale approaches the relaxation timescale. The magnitude of these peaks varies systematically with Mach number and wall temperature. At Mach~6 with $T_w = 1200$~K (panel~(a)), moderate peaks emerge, whilst elevated wall temperature (panel~(b)) amplifies them significantly. At Mach~8 (panels~(c) and~(d)), peak magnitudes are substantially reduced due to faster viscous timescales, with cold wall conditions (panel~(c)) showing the most suppression. These results contrast with the laminar findings of Chapter~4, demonstrating that the turbulent regime exhibits different dynamics where the balance between viscous timescale compression and enhanced vibrational energy reshapes the spatial confinement and intensity of thermal nonequilibrium.


\begin{figure}[t]
    \centering
    \fontsize{9}{10}\selectfont
    \includegraphics[width=0.80\textwidth]{resources/figures/breakdown/result_iles_Damkohler_number.pdf}
    \caption[Wall-normal profiles of Damk\"{o}hler number across streamwise stations for TNE formulations at (a) Mach 6, $T_w=1200$ K; (b) Mach 6, $T_w=1800$ K; (c) Mach 8, $T_w=1200$ K; and (d) Mach 8, $T_w=1800$ K.]{Wall-normal profiles of Damk\"{o}hler number defined as the ratio of viscous timescale to vibrational relaxation timescale across five streamwise stations (laminar, resonance, secondary instability, transition, and turbulent regions) for TNE formulations showing (a) Mach 6, $T_w=1200$ K; (b) Mach 6, $T_w=1800$ K; (c) Mach 8, $T_w=1200$ K; and (d) Mach 8, $T_w=1800$ K. Colours denote streamwise position, with darker shades indicating downstream evolution.}
    \label{figure:damkohlernumber}
\end{figure}


% \begin{figure}[h]
%     \centering
%     \fontsize{9}{10}\selectfont
%     \includegraphics[width=0.80\textwidth]{resources/figures/breakdown/result_dns_thermal_rhoT_rhoTv.pdf}
%     \caption[Streamwise evolution of spatially temporally averaged translational--rotational ($C_{q,tr}$, black) and vibrational ($C_{q,v}$, red, scaled $\times 10$) heat flux coefficients for the Mach~6, $T_w = 1800$ K case. Results are shown for TNE formulations, (a) M6Tw12, (b) M6Tw18, (c) M8Tw12, (d) M8Tw18]%
%     {Results are shown for \gls{TNE} formulations, (a) M6Tw12, (b) M6Tw18, (c) M8Tw12, (d) M8Tw18.}
%     \label{figure:rhoT_rhoTv}
% \end{figure}




% The streamwise evolution of the Damk\"{o}hler number at five characteristic stations spanning the laminar region, resonance zone, secondary instability region, transition zone, and fully turbulent region is presented in \ref{figure:damkohlernumber}. The Damk\"{o}hler number, defined as the ratio of the viscous timescale $\delta^{*}/u_{\tau}$ to the vibrational relaxation timescale, quantifies the relative importance of instability-driven energy redistribution compared to molecular relaxation processes. At Mach~6 with $T_w = 1200$~K (panel~(a)), the Damk\"{o}hler number remains relatively low in the laminar region, remaining confined to a thin wall-adjacent layer. However, as the flow progresses through the secondary instability region, a pronounced peak emerges and shifts away from the wall, indicating that secondary instabilities drive hydrodynamic acceleration throughout a broader region of the boundary layer. Elevated wall temperature at Mach~6 (panel~(b)) amplifies these peaks substantially and extends their wall-normal reach, reflecting enhanced vibrational excitation across a thicker layer. At Mach~8, the distributions show fundamentally different behaviour: at cold wall conditions (panel~(c)) the peaks are significantly reduced and remain tightly confined near the wall, whilst at hot wall conditions (panel~(d)) the magnitudes recover but still maintain a narrower wall-normal extent compared to Mach~6. These distributions confirm the laminar stability trends from Chapter~4, and the consistent inward migration of elevated Damk\"{o}hler number profiles towards the wall during transition demonstrates that vibrational excitation becomes increasingly wall-confined as secondary instabilities develop, with the intensity and spatial distribution of this confinement controlled by the interplay between Mach number and wall temperature effects on molecular relaxation timescales.

% The left panels of \Cref{figure:rhoT_rhoTv} show mean temperature profiles $\tilde{T}_i / T_\infty$ plotted against $y^*$ on a logarithmic scale. The CPG model exhibits the steepest temperature gradient near the wall, followed by TEQ, whilst TNE displays a more gradual temperature rise into the freestream. This difference in mean temperature structure directly reflects the altered thermodynamic response of the fluid when vibrational energy remains decoupled from translational energy. In the CPG and TEQ cases, the full sensible enthalpy is distributed between translational kinetic energy and internal energy, both of which respond immediately to molecular collisions and heat transfer. Conversely, under TNE conditions, a substantial fraction of thermal energy is sequestered in vibrational degrees of freedom, which exchange energy on much longer timescales. This partitioning acts to reduce the effective temperature rise in the translational modes, thereby lowering local Mach numbers and reducing the growth rate of temperature-sensitive instabilities, the Mack modes.

% The right panels present turbulent temperature fluctuations $\sqrt{\overline{T'^2}} / \tilde{T}_\infty$, which provide the most direct measure of transition onset and development. At small $y^*$ (inner layer), all three models show negligible fluctuations, corresponding to the laminar state. However, as $y^*$ increases towards the transition region, the models diverge markedly. The TEQ case exhibits the earliest and most vigorous onset of temperature fluctuations, with a pronounced peak that rises steeply and reaches maximum amplitude in the intermediate region. CPG follows closely behind, showing comparable peak magnitude but slightly delayed emergence. Most significantly, TNE exhibits substantially suppressed temperature fluctuations throughout the transition zone, with peak amplitudes approximately $30$--$40\,\%$ lower than TEQ and a much broader, flatter plateau that extends further into larger $y^*$ values. This suppression is not merely a quantitative difference; it represents a fundamental alteration of the transition process itself.

Finally, the vibrational temperature deviation profiles presented in \ref{figure:iles_tmtv_profiles} across both linear and semi-logarithmic coordinates reveal the wall-normal progression of nonequilibrium through the streamwise stations. Consistent with the Damköhler number analysis, the vibrational temperature deficit progresses towards the wall as the flow transitions through the secondary instability region. In the laminar station, the vibrational temperature deviation remains distributed throughout the boundary layer, but as secondary instabilities develop, the profiles sharpen and concentrate their maximum gradient in the wall-adjacent region. This progressive wall-confinement reflects the same timescale balance discussed previously, nonequilibrium effects intensify where viscous processes and molecular relaxation become comparable. The semi-logarithmic coordinates (right column) reveal the inner-layer structure more clearly, exposing pronounced features in the buffer layer ($y^+ \sim 5$--$30$) and logarithmic layer ($y^+ \sim 30$--$300$) across the thermal conditions. At Mach~6, particularly at elevated wall temperature, the vibrational temperature deviation exhibits a secondary peak in the buffer layer during transition, indicating localised energy redistribution driven by the intense shear and secondary instability dynamics. At Mach~8, these buffer-layer features are substantially suppressed, consistent with the reduced Damköhler number magnitudes, demonstrating that the faster viscous timescale inhibits vibrational excitation regardless of wall temperature effects. The semi-log visualisation thus confirms that the inner-layer dynamics controlling thermal nonequilibrium are fundamentally reshaped by the competing influences of Mach number and wall temperature on the viscous and relaxation timescales.

\begin{figure}[t]
    \centering
    \fontsize{9}{10}\selectfont
    \includegraphics[width=0.80\textwidth]{resources/figures/breakdown/result_iles_thermal_TmTv.pdf}
    \caption[Wall-normal profiles of vibrational temperature deviation normalised by freestream conditions across streamwise stations for TNE formulations at (a) Mach 6, $T_w=1200$ K; (b) Mach 6, $T_w=1800$ K; (c) Mach 8, $T_w=1200$ K; and (d) Mach 8, $T_w=1800$ K.]{Wall-normal profiles of vibrational temperature deviation $(T_v - T)/T_\infty$ across five streamwise stations (laminar, resonance, secondary instability, transition, and turbulent regions) for TNE formulations showing (a) Mach 6, $T_w=1200$ K; (b) Mach 6, $T_w=1800$ K; (c) Mach 8, $T_w=1200$ K; and (d) Mach 8, $T_w=1800$ K. Left column presents linear wall-normal coordinates $y/\delta^{*}$, whilst right column presents semi-logarithmic coordinates to reveal inner-layer structure. Colours denote streamwise position, with darker shades indicating downstream evolution.}
    \label{figure:iles_tmtv_profiles}
\end{figure}


% The physical origin of this suppression lies in the energy redistribution mechanism of nonequilibrium. In TNE flow, the Landau--Teller relaxation process governs vibrational energy exchange. When turbulent disturbances attempt to amplify through nonlinear interactions, the energy available for amplification is partitioned between translational and vibrational modes according to the instantaneous relaxation time. Since vibrational energy relaxes on timescales of $\mathcal{O}(10^{-6}$--$10^{-5})$ seconds---comparable to or longer than the local eddy turnover time---the translational temperature cannot respond as quickly to disturbance-induced pressure and shear fluctuations. This decoupling effectively damps the amplification of turbulent structures in their early growth phase, explaining the delayed and attenuated rise of $\sqrt{\overline{T'^2}}$ under TNE conditions.

% The consequence of this mechanism is a systematic ordering of transition location: TEQ transitions earliest, followed by CPG at an intermediate streamwise position, whilst TNE exhibits the most delayed transition. This ordering has profound implications for heat transfer prediction and stability analysis, since the transition location determines the extent of laminar and turbulent domains and hence the integrated heat load on the vehicle surface. Importantly, the TNE effect persists across all examined wall temperatures and Mach numbers, indicating a robust and thermodynamically fundamental mechanism rather than a case-specific artefact.



\subsection{Second-order statistics}

% \subsubsection{Heat flux and Reynolds stresses}
% Similar story of ILES domain getting this correct

% \subsubsection{Second-order statistics}




% \newpage

% \section{Chapter Summary}
% Key points to drive home, TEQ transitions faster than CPG and TNE, with TNE showing significant delay in transition. Grid study shows that DNS can caputure the transition location compared to ILES. ILES was used to understand the behaviour of Mach numbers and wall temperatures for a TNE case. 

\section*{Chapter summary}

Direct numerical simulation of hypersonic boundary layer transition is performed across calorically perfect gas (CPG), thermal equilibrium (TEQ), and thermal nonequilibrium (TNE) formulations at Mach~6 and Mach~8 with wall temperatures of 1200~K and 1800~K. Second-mode instabilities are excited via controlled forcing, revealing a robust transition hierarchy: TEQ transitions earliest, CPG at intermediate streamwise location, and TNE exhibits the most delayed transition by approximately 150 displacement thicknesses. This ordering, driven by energy redistribution through the Landau--Teller relaxation process, persists across all Mach numbers and wall temperatures, with elevated wall temperature and higher Mach number further stabilising transition by reducing effective Mach numbers and suppressing secondary instability growth. During the laminar and early transitional stages, TNE effects are modest; however, as breakdown proceeds, vibrational energy sequestration increasingly damps turbulent structure amplification, elevating secondary instability thresholds and attenuating streak development compared to CPG and TEQ. This is reflected in suppressed overshoot of skin friction and Stanton number in TNE compared to the pronounced overshoot seen in TEQ. Notably, once turbulence is fully developed, velocity Reynolds stresses collapse across all three models, whilst temperature fluctuations remain markedly elevated in TNE, demonstrating that vibrational nonequilibrium effects are primarily important during the transition process itself and exert lasting influence on thermal fluctuations in the turbulent regime. These findings underscore the necessity of accounting for thermal nonequilibrium when predicting hypersonic transition and surface heating across the parameter space relevant to hypersonic vehicle design.

